% =============================================================================
% QUESTIONÁRIO WOCAT SOBRE TECNOLOGIAS DE MANEJO SUSTENTÁVEL DA TERRA
% Tradução Integral para Português Brasileiro (2019)
% =============================================================================
% Traduzido no âmbito do Artigo 3 da tese de doutorado de
% Catuxe Varjão de Santana Oliveira — PPGPI/UFS
% =============================================================================

\documentclass[12pt,a4paper,brazil]{article}

\usepackage[utf8]{inputenc}
\usepackage[T1]{fontenc}
\usepackage[brazil]{babel}
\usepackage{geometry}
\usepackage{setspace}
\usepackage{indentfirst}
\usepackage{enumitem}
\usepackage{booktabs}
\usepackage{array}
\usepackage{longtable}
\usepackage{graphicx}
\usepackage{calc}
\usepackage{amssymb}
\usepackage{hyperref}
\usepackage{titlesec}
\usepackage{fancyhdr}

% Definição necessária para conteúdo gerado pelo Pandoc
\providecommand{\tightlist}{\setlength{\itemsep}{0pt}\setlength{\parskip}{0pt}}

\geometry{a4paper, left=3cm, right=2cm, top=3cm, bottom=2cm}
\onehalfspacing

% Cabeçalho e rodapé
\pagestyle{fancy}
\fancyhf{}
\fancyhead[L]{\small Questionário WOCAT --- Tecnologias de MST}
\fancyhead[R]{\small Tradução PT-BR (2019)}
\fancyfoot[C]{\thepage}
\setlength{\headheight}{14.5pt}

\hypersetup{
    colorlinks=true,
    linkcolor=blue,
    urlcolor=blue,
    pdfauthor={Catuxe Varjão de Santana Oliveira},
    pdftitle={Questionário WOCAT - Tecnologias de MST - Tradução PT-BR},
    pdfsubject={Manejo Sustentável da Terra}
}

\begin{document}

% Capa
\begin{titlepage}
\centering
\vspace*{3cm}

{\Large\bfseries QUESTIONÁRIO WOCAT}\\[0.5cm]
{\Large\bfseries SOBRE TECNOLOGIAS DE}\\[0.5cm]
{\Large\bfseries MANEJO SUSTENTÁVEL DA TERRA (MST)}\\[1.5cm]

{\large Tradução Integral para Português Brasileiro}\\[0.3cm]
{\large Versão 2019}\\[3cm]

{\normalsize Traduzido no âmbito da tese de doutorado:}\\[0.3cm]
{\normalsize\itshape ``Gestão da Propriedade Intelectual e Conservação do}\\
{\normalsize\itshape Conhecimento Tradicional Quilombola: Adaptação Transcultural}\\
{\normalsize\itshape do Questionário WOCAT-SLM e Construção do Índice de}\\
{\normalsize\itshape Sustentabilidade Bioculural (ISB)''}\\[2cm]

{\normalsize Catuxe Varjão de Santana Oliveira}\\[0.3cm]
{\normalsize Programa de Pós-Graduação em Propriedade Intelectual (PPGPI)}\\[0.3cm]
{\normalsize Universidade Federal de Sergipe (UFS)}\\[0.3cm]
{\normalsize 2026}

\vfill
{\footnotesize Fonte original: WOCAT --- World Overview of Conservation Approaches and Technologies\\
\url{https://www.wocat.net}}

\end{titlepage}

\tableofcontents
\newpage

% Conteúdo do questionário
\subsection{Introdução ao questionário}\label{introduuxe7uxe3o-ao-questionuxe1rio}

\subsubsection{Sobre a documentação WOCAT de práticas de MST}\label{sobre-a-documentauxe7uxe3o-wocat-de-pruxe1ticas-de-mst}

O WOCAT fornece ferramentas padronizadas, dirigidas pelo usuário, de acesso aberto e uso global para documentação e avaliação de práticas de Manejo Sustentável da Terra (MST). \textbf{MST}, no contexto do WOCAT, define-se como o uso sustentável dos recursos da terra, incluindo solos, água, vegetação e animais. O WOCAT concentra-se em esforços para prevenir e reduzir a degradação da terra e restaurar terras degradadas por meio de \textbf{tecnologias de manejo da terra} e \textbf{abordagens para implementá-las}. Todas as práticas podem ser consideradas, sejam elas indígenas, recentemente introduzidas por projetos ou inovações recentes de usuários da terra. Todas as informações documentadas por meio de questionários WOCAT são disponibilizadas no Banco de Dados Global de MST e podem ser usadas para disseminar conhecimento sobre MST e melhorar a tomada de decisões.

\subsubsection{Tecnologia ou Abordagem?}\label{tecnologia-ou-abordagem}

Existem dois questionários separados: um para Tecnologias e outro para Abordagens. Juntos, fornecem o panorama completo de uma prática de MST. Idealmente, preencha primeiro o questionário sobre Tecnologias de MST, seguido pelo questionário sobre Abordagens de MST.

\begin{itemize}
\item
  Uma \textbf{Tecnologia de MST} é uma prática física que controla a degradação da terra e aumenta a produtividade e/ou outros serviços ecossistêmicos. Uma Tecnologia consiste em uma ou várias medidas, como medidas agronômicas, vegetativas, estruturais e de manejo.
\item
  Uma \textbf{Abordagem de MST} define os meios empregados para implementar uma ou mais Tecnologias de MST. Inclui suporte técnico e material, bem como o envolvimento e papéis de diferentes partes interessadas.
\end{itemize}

Uma Abordagem deve sempre estar vinculada a uma ou mais Tecnologias.

\subsubsection{Como documentar e revisar dados WOCAT}\label{como-documentar-e-revisar-dados-wocat}

\begin{enumerate}
\def\labelenumi{\arabic{enumi}.}
\tightlist
\item
  Familiarize-se com o questionário em papel. Leia as instruções, explicações, definições e exemplos.
\item
  Comece a preencher o questionário com base em seu conhecimento e documentos existentes.
\item
  Identifique usuários da terra e outras pessoas-recurso com conhecimento aprofundado da Tecnologia/Abordagem de MST.
\item
  Colete dados em campo. Reúna informações por meio de entrevistas com usuário(s) da terra e pessoas-recurso. Faça medições, tire fotos e elabore desenhos técnicos.
\item
  Insira as informações compiladas no Banco de Dados Global de MST em \url{https://qcat.wocat.net}.
\item
  O Banco de Dados Global orientará sobre como editar e submeter seus dados para revisão.
\end{enumerate}

\subsubsection{Observações gerais}\label{observauxe7uxf5es-gerais}

\begin{itemize}
\tightlist
\item
  \textbf{Responda a todas as perguntas.} Se dados precisos não estiverem disponíveis, forneça a melhor estimativa com base em seu julgamento profissional. Se determinadas perguntas não forem aplicáveis ou relevantes, indique ``n/a''.
\item
  Perguntas com o ícone {[Consulta]} devem ser respondidas em consulta com usuários da terra.
\item
  Perguntas com o ícone {[Campo]} requerem medições ou observações em campo.
\item
  Círculos ($\circ$) indicam pergunta de seleção única. Marque apenas uma resposta.
\item
  Quadrados ($\Box$) permitem selecionar várias respostas.
\item
  \textbf{Utilize documentos existentes e consulte outros especialistas em MST e usuários da terra para melhorar a qualidade dos dados.}
\item
  Preencha um questionário separado para cada Tecnologia e para cada Abordagem.
\end{itemize}

\textbf{Fonte:} WOCAT --- \emph{World Overview of Conservation Approaches and Technologies}, Centre for Development and Environment (CDE), Universidade de Berna, Suíça. Versão 2019.

\newpage

\begin{center}\rule{0.5\linewidth}{0.5pt}\end{center}

\subsection{1. Informações Gerais}\label{informauxe7uxf5es-gerais}

\subsubsection{1.1 Nome da Tecnologia de MST (doravante referida como ``Tecnologia'')}\label{nome-da-tecnologia-de-mst-doravante-referida-como-tecnologia}

{
\begin{longtable}[]{@{}ll@{}}
\toprule\noalign{}
Campo & Resposta \\
\midrule\noalign{}
\endhead
\bottomrule\noalign{}
\endlastfoot
Nome & \_\_\_\_\_\_\_\_\_\_\_\_\_\_\_\_\_\_\_\_\_\_\_\_\_\_\_\_\_\_\_\_\_\_\_\_\_\_\_\_ \\
Nome localmente utilizado & \_\_\_\_\_\_\_\_\_\_\_\_\_\_\_\_\_\_\_\_\_\_\_\_\_\_\_\_\_\_\_\_\_\_\_\_\_\_\_\_ \\
País & \_\_\_\_\_\_\_\_\_\_\_\_\_\_\_\_\_\_\_\_\_\_\_\_\_\_\_\_\_\_\_\_\_\_\_\_\_\_\_\_ \\
\end{longtable}
}

\subsubsection{1.2 Dados de contato das pessoas-recurso e instituições envolvidas na avaliação e documentação da Tecnologia}\label{dados-de-contato-das-pessoas-recurso-e-instituiuxe7uxf5es-envolvidas-na-avaliauxe7uxe3o-e-documentauxe7uxe3o-da-tecnologia}

\subsubsection{Compilador(a)}\label{compiladora}

\emph{A pessoa que conduziu as entrevistas, compilou as informações e preencheu o questionário.}

{
\begin{longtable}[]{@{}ll@{}}
\toprule\noalign{}
Campo & Resposta \\
\midrule\noalign{}
\endhead
\bottomrule\noalign{}
\endlastfoot
Sobrenome & \_\_\_\_\_\_\_\_\_\_\_\_\_\_\_\_\_\_\_\_\_\_\_\_\_\_\_\_\_\_\_\_\_\_\_\_\_\_\_\_ \\
Nome(s) & \_\_\_\_\_\_\_\_\_\_\_\_\_\_\_\_\_\_\_\_\_\_\_\_\_\_\_\_\_\_\_\_\_\_\_\_\_\_\_\_ \\
Gênero & $\circ$ Sra. $\circ$ Sr. \\
Nome da instituição & \_\_\_\_\_\_\_\_\_\_\_\_\_\_\_\_\_\_\_\_\_\_\_\_\_\_\_\_\_\_\_\_\_\_\_\_\_\_\_\_ \\
País & \_\_\_\_\_\_\_\_\_\_\_\_\_\_\_\_\_\_\_\_\_\_\_\_\_\_\_\_\_\_\_\_\_\_\_\_\_\_\_\_ \\
Telefone & \_\_\_\_\_\_\_\_\_\_\_\_\_\_\_\_\_\_\_\_\_\_\_\_\_\_\_\_\_\_\_\_\_\_\_\_\_\_\_\_ \\
E-mail & \_\_\_\_\_\_\_\_\_\_\_\_\_\_\_\_\_\_\_\_\_\_\_\_\_\_\_\_\_\_\_\_\_\_\_\_\_\_\_\_ \\
\end{longtable}
}

\subsubsection{Pessoa(s)-recurso principal(is)}\label{pessoas-recurso-principalis}

\emph{Pessoa(s) que forneceram a maioria das informações documentadas neste questionário. Podem ser usuários da terra, especialistas em MST (p.ex. assessores técnicos, pesquisadores) ou outras pessoas.}

\textbf{Pessoa-recurso 1:}

{
\begin{longtable}[]{@{}
  >{\raggedright\arraybackslash}p{(\linewidth - 2\tabcolsep) * \real{0.4118}}
  >{\raggedright\arraybackslash}p{(\linewidth - 2\tabcolsep) * \real{0.5882}}@{}}
\toprule\noalign{}
\begin{minipage}[b]{\linewidth}\raggedright
Campo
\end{minipage} & \begin{minipage}[b]{\linewidth}\raggedright
Resposta
\end{minipage} \\
\midrule\noalign{}
\endhead
\bottomrule\noalign{}
\endlastfoot
Tipo & $\circ$ Usuário da terra $\circ$ Especialista/assessor técnico $\circ$ Co-compilador $\circ$ Outro: \_\_\_\_\_\_\_\_ \\
Registro WOCAT & $\circ$ Usuário registrado $\circ$ Não registrado \\
Sobrenome & \_\_\_\_\_\_\_\_\_\_\_\_\_\_\_\_\_\_\_\_\_\_\_\_\_\_\_\_\_\_\_\_\_\_\_\_\_\_\_\_ \\
Nome(s) & \_\_\_\_\_\_\_\_\_\_\_\_\_\_\_\_\_\_\_\_\_\_\_\_\_\_\_\_\_\_\_\_\_\_\_\_\_\_\_\_ \\
Instituição & \_\_\_\_\_\_\_\_\_\_\_\_\_\_\_\_\_\_\_\_\_\_\_\_\_\_\_\_\_\_\_\_\_\_\_\_\_\_\_\_ \\
País & \_\_\_\_\_\_\_\_\_\_\_\_\_\_\_\_\_\_\_\_\_\_\_\_\_\_\_\_\_\_\_\_\_\_\_\_\_\_\_\_ \\
\end{longtable}
}

\emph{Repita para pessoas-recurso adicionais (2, 3, 4), conforme necessário.}

\textbf{Nome da(s) instituição(ões)} que facilitaram a documentação/avaliação da Tecnologia: \_\_\_\_\_\_\_\_\_\_\_\_\_\_\_\_\_\_\_\_\_\_\_\_\_\_\_\_\_\_\_\_\_\_\_\_\_\_\_\_

\textbf{Nome do projeto} que facilitou a documentação/avaliação da Tecnologia: \_\_\_\_\_\_\_\_\_\_\_\_\_\_\_\_\_\_\_\_\_\_\_\_\_\_\_\_\_\_\_\_\_\_\_\_\_\_\_\_

\subsubsection{1.3 Condições relativas ao uso de dados documentados via WOCAT}\label{condiuxe7uxf5es-relativas-ao-uso-de-dados-documentados-via-wocat}

O compilador e a(s) pessoa(s)-recurso aceitam as condições relativas ao uso de dados documentados via WOCAT?

$\circ$ Sim $\circ$ Não

\emph{Nota: Caso não aceite as condições, não será possível inserir e editar dados no banco de dados WOCAT.}

\textbf{Condições:}

\begin{itemize}
\tightlist
\item
  Os dados capturados serão inseridos, editados e armazenados no banco de dados online do WOCAT pelo compilador. A responsabilidade geral pela compilação e qualidade dos dados é do compilador.
\item
  Os dados armazenados no banco de dados WOCAT são de acesso aberto.
\item
  Os dados são disponibilizados sob a licença \emph{Creative Commons Atribuição-NãoComercial-CompartilhaIgual 3.0 Não Adaptada} (CC BY-NC-SA 3.0).
\end{itemize}

\subsubsection{1.4 Declaração sobre a sustentabilidade da Tecnologia descrita}\label{declarauxe7uxe3o-sobre-a-sustentabilidade-da-tecnologia-descrita}

A Tecnologia possui efeitos adversos sobre a degradação da terra, de modo que não pode ser declarada uma tecnologia de manejo \emph{sustentável} da terra?

$\circ$ Sim $\circ$ Não

Comentários: \_\_\_\_\_\_\_\_\_\_\_\_\_\_\_\_\_\_\_\_\_\_\_\_\_\_\_\_\_\_\_\_\_\_\_\_\_\_\_\_

\subsubsection{1.5 Referência a Questionário(s) sobre Abordagens de MST (documentados usando WOCAT)}\label{referuxeancia-a-questionuxe1rios-sobre-abordagens-de-mst-documentados-usando-wocat}

{
\begin{longtable}[]{@{}ll@{}}
\toprule\noalign{}
Campo & Resposta \\
\midrule\noalign{}
\endhead
\bottomrule\noalign{}
\endlastfoot
Nome da Abordagem de MST & \_\_\_\_\_\_\_\_\_\_\_\_\_\_\_\_\_\_\_\_\_\_\_\_\_\_\_\_\_\_\_\_\_\_\_\_\_\_\_\_ \\
Compilador & \_\_\_\_\_\_\_\_\_\_\_\_\_\_\_\_\_\_\_\_\_\_\_\_\_\_\_\_\_\_\_\_\_\_\_\_\_\_\_\_ \\
\end{longtable}
}

\newpage

\begin{center}\rule{0.5\linewidth}{0.5pt}\end{center}

\subsection{2. Descrição de uma Tecnologia de MST}\label{descriuxe7uxe3o-de-uma-tecnologia-de-mst}

\emph{Uma Tecnologia de MST é uma prática aplicada no campo que controla a degradação da terra e/ou aumenta a produtividade. Este questionário foi projetado para documentar uma única Tecnologia de MST e não pode ser usado para avaliar uma propriedade inteira.}

\emph{Uma Tecnologia de MST pode consistir em uma ou várias medidas de MST (agronômicas, vegetativas, estruturais e de manejo); p.ex. terraços combinados com faixas de gramíneas e aração em nível.}

\subsubsection{2.1 Descrição curta da Tecnologia}\label{descriuxe7uxe3o-curta-da-tecnologia}

\emph{Resuma a Tecnologia em 1--2 frases. Certifique-se de que a descrição curta é precisa e contém palavras-chave relevantes.}

\begin{center}\rule{0.5\linewidth}{0.5pt}\end{center}

\begin{center}\rule{0.5\linewidth}{0.5pt}\end{center}

\subsubsection{2.2 Descrição detalhada da Tecnologia {[Consulta]} {[Campo]}}\label{descriuxe7uxe3o-detalhada-da-tecnologia}

\emph{A descrição detalhada deve apresentar um panorama conciso mas abrangente da Tecnologia para leitores externos. Deve abordar questões-chave como:}

\begin{enumerate}
\def\labelenumi{\arabic{enumi}.}
\tightlist
\item
  \emph{Quais são as principais características/elementos da Tecnologia (incluindo especificações técnicas)?}
\item
  \emph{Onde a Tecnologia é aplicada (ambiente natural e humano)?}
\item
  \emph{Quais são os propósitos/funções da Tecnologia?}
\item
  \emph{Quais as principais atividades/insumos necessários para implantar/manter a Tecnologia?}
\item
  \emph{Quais são os benefícios/impactos da Tecnologia?}
\item
  \emph{O que os usuários da terra apreciam/não apreciam na Tecnologia?}
\end{enumerate}

\emph{Extensão ideal: 2.500--3.000 caracteres; máximo absoluto: 3.500 caracteres.}

\begin{center}\rule{0.5\linewidth}{0.5pt}\end{center}

\begin{center}\rule{0.5\linewidth}{0.5pt}\end{center}

\begin{center}\rule{0.5\linewidth}{0.5pt}\end{center}

\begin{center}\rule{0.5\linewidth}{0.5pt}\end{center}

\begin{center}\rule{0.5\linewidth}{0.5pt}\end{center}

\subsubsection{2.3 Fotografias da Tecnologia {[Campo]}}\label{fotografias-da-tecnologia}

\emph{Forneça fotos mostrando visão geral e detalhes da Tecnologia. Ao menos dois arquivos digitais (JPG, PNG, GIF) de alta qualidade/resolução. Uma legenda explicativa é obrigatória para cada foto. Fotos devem ilustrar o desenho técnico (4.1) e, quando apropriado, mostrar a situação antes/depois ou com/sem medidas de MST.}

{
\begin{longtable}[]{@{}lllll@{}}
\toprule\noalign{}
Nome do arquivo & Legenda/explicação & Data & Local & Fotógrafo \\
\midrule\noalign{}
\endhead
\bottomrule\noalign{}
\endlastfoot
& & & & \\
& & & & \\
& & & & \\
\end{longtable}
}

\subsubsection{2.4 Vídeos da Tecnologia}\label{vuxeddeos-da-tecnologia}

\emph{Se disponíveis, carregue em plataforma pública (p.ex. vimeo.com, youtube.com) e indique link e descrição curta.}

{
\begin{longtable}[]{@{}lllll@{}}
\toprule\noalign{}
Link & Descrição curta & Data & Local & Videógrafo \\
\midrule\noalign{}
\endhead
\bottomrule\noalign{}
\endlastfoot
& & & & \\
\end{longtable}
}

\subsubsection{2.5 País/região/locais onde a Tecnologia foi aplicada e que são cobertos por esta avaliação}\label{pauxedsregiuxe3olocais-onde-a-tecnologia-foi-aplicada-e-que-suxe3o-cobertos-por-esta-avaliauxe7uxe3o}

\emph{Restrinja as informações apenas aos locais que foram avaliados/analisados no processo de documentação.}

{
\begin{longtable}[]{@{}
  >{\raggedright\arraybackslash}p{(\linewidth - 2\tabcolsep) * \real{0.4118}}
  >{\raggedright\arraybackslash}p{(\linewidth - 2\tabcolsep) * \real{0.5882}}@{}}
\toprule\noalign{}
\begin{minipage}[b]{\linewidth}\raggedright
Campo
\end{minipage} & \begin{minipage}[b]{\linewidth}\raggedright
Resposta
\end{minipage} \\
\midrule\noalign{}
\endhead
\bottomrule\noalign{}
\endlastfoot
País & \_\_\_\_\_\_\_\_\_\_\_\_\_\_\_\_\_\_\_\_\_\_\_\_\_\_\_\_\_\_\_\_\_\_\_\_\_\_\_\_ \\
Região/Estado/Província & \_\_\_\_\_\_\_\_\_\_\_\_\_\_\_\_\_\_\_\_\_\_\_\_\_\_\_\_\_\_\_\_\_\_\_\_\_\_\_\_ \\
Especificação adicional do local & \_\_\_\_\_\_\_\_\_\_\_\_\_\_\_\_\_\_\_\_\_\_\_\_\_\_\_\_\_\_\_\_\_\_\_\_\_\_\_\_ \\
\end{longtable}
}

Número de locais considerados/analisados na documentação desta Tecnologia:

$\circ$ Local único $\circ$ 2--10 locais $\circ$ 10--100 locais $\circ$ 100--1.000 locais $\circ$ \textgreater{} 1.000 locais

\textbf{Coordenadas georreferenciadas (graus decimais):}

{
\begin{longtable}[]{@{}lll@{}}
\toprule\noalign{}
Nome do local / usuário & Latitude & Longitude \\
\midrule\noalign{}
\endhead
\bottomrule\noalign{}
\endlastfoot
& & \\
& & \\
\end{longtable}
}

Especifique a distribuição da Tecnologia:

$\circ$ Distribuída uniformemente sobre uma área $\circ$ Aplicada em pontos específicos/concentrada em área pequena

Se distribuída uniformemente, área coberta (km²): \_\_\_\_\_\_\_\_

O(s) local(is) da Tecnologia está(ão) em área permanentemente protegida? $\circ$ Sim $\circ$ Não

\subsubsection{2.6 Data de implementação}\label{data-de-implementauxe7uxe3o}

Ano de implementação: \_\_\_\_\_\_\_\_

Se o ano exato não for conhecido, indique a data aproximada:

$\circ$ Menos de 10 anos (recente) $\circ$ 10--50 anos $\circ$ Mais de 50 anos (tradicional)

\subsubsection{2.7 Introdução da Tecnologia {[Campo]}}\label{introduuxe7uxe3o-da-tecnologia}

\emph{Várias respostas possíveis.}

Especifique como a Tecnologia foi introduzida:

\begin{itemize}
\tightlist
\item[$\square$]
  Como parte de um sistema tradicional
\item[$\square$]
  Por inovação recente de usuários da terra
\item[$\square$]
  Durante experimentos/pesquisa
\item[$\square$]
  Por projetos/intervenções externas
\item[$\square$]
  Outro (especifique): \_\_\_\_\_\_\_\_\_\_\_\_\_\_\_\_\_\_\_\_\_\_\_\_\_\_\_\_\_\_\_\_\_\_\_\_\_\_\_\_
\end{itemize}

Comentários (tipo de projeto, etc.): \_\_\_\_\_\_\_\_\_\_\_\_\_\_\_\_\_\_\_\_\_\_\_\_\_\_\_\_\_\_\_\_\_\_\_\_\_\_\_\_

\newpage

\begin{center}\rule{0.5\linewidth}{0.5pt}\end{center}

\subsection{3. Classificação da Tecnologia de MST}\label{classificauxe7uxe3o-da-tecnologia-de-mst}

\subsubsection{3.1 Propósito(s) principal(is) da Tecnologia {[Campo]}}\label{propuxf3sitos-principalis-da-tecnologia}

\emph{Várias respostas possíveis. Máximo 5 respostas.}

\begin{itemize}
\tightlist
\item[$\square$]
  Melhorar a produção (culturas, forragem, madeira/fibra, água, energia)
\item[$\square$]
  Prevenir (evitar), reduzir a degradação da terra; restaurar/reabilitar a terra (solo, água, vegetação)
\item[$\square$]
  Conservar ecossistemas
\item[$\square$]
  Preservar/melhorar a biodiversidade
\item[$\square$]
  Criar impacto econômico benéfico (p.ex. aumentar renda/oportunidades de emprego)
\item[$\square$]
  Criar impacto social benéfico (p.ex. reduzir conflitos por recursos naturais, apoiar grupos marginalizados)
\item[$\square$]
  Reduzir risco de desastres (p.ex. secas, enchentes, deslizamentos)
\item[$\square$]
  Adaptar-se a mudanças/extremos climáticos e seus impactos (p.ex. resiliência a secas, tempestades)
\item[$\square$]
  Mitigar mudanças climáticas e seus impactos (p.ex. sequestro de carbono)
\item[$\square$]
  Outro propósito (especifique): \_\_\_\_\_\_\_\_\_\_\_\_\_\_\_\_\_\_\_\_\_\_\_\_\_\_\_\_\_\_\_\_\_\_\_\_\_\_\_\_
\end{itemize}

\subsubsection{3.2 Tipo(s) atual(is) de uso da terra onde a Tecnologia é aplicada {[Campo]}}\label{tipos-atualis-de-uso-da-terra-onde-a-tecnologia-uxe9-aplicada}

O uso da terra é misto na mesma unidade (conforme definições ICRAF)?

$\circ$ Sim $\circ$ Não

Se sim, especifique o uso misto em sistema agroflorestal:

$\circ$ Agrossilvicultura (p.ex. lavoura e árvores) $\circ$ Agrossilvipastoril (lavoura + pastagem/animais + árvores) $\circ$ Silvipastoril (árvores e pastagem/animais)

\subsubsection{Tipos de uso da terra e subcategorias}\label{tipos-de-uso-da-terra-e-subcategorias}

{
\begin{longtable}[]{@{}
  >{\raggedright\arraybackslash}p{(\linewidth - 4\tabcolsep) * \real{0.2955}}
  >{\raggedright\arraybackslash}p{(\linewidth - 4\tabcolsep) * \real{0.3409}}
  >{\raggedright\arraybackslash}p{(\linewidth - 4\tabcolsep) * \real{0.3636}}@{}}
\toprule\noalign{}
\begin{minipage}[b]{\linewidth}\raggedright
Tipo de uso
\end{minipage} & \begin{minipage}[b]{\linewidth}\raggedright
Subcategorias
\end{minipage} & \begin{minipage}[b]{\linewidth}\raggedright
Especificações
\end{minipage} \\
\midrule\noalign{}
\endhead
\bottomrule\noalign{}
\endlastfoot
$\Box$ \textbf{Terra de cultivo} & $\Box$ Cultivo anual; $\Box$ Cultivo perene; $\Box$ Cultivo arbóreo/arbustivo; $\Box$ Outro & Culturas: \_\_\_\_\_\_\_\_ Estações de crescimento/ano: $\circ$ 1 $\circ$ 2 $\circ$ 3 --- Rotação de culturas? $\circ$ Sim $\circ$ Não --- Consórcio? $\circ$ Sim $\circ$ Não \\
$\Box$ \textbf{Terra de pastagem} & \textbf{Extensivo:} $\Box$ Nomadismo; $\Box$ Seminomadismo; $\Box$ Transumância; $\Box$ Ranching --- \textbf{Intensivo:} $\Box$ Corte-e-carrega/confinamento; $\Box$ Pastagem melhorada & Tipo de animal: \_\_\_\_\_\_\_\_ Integração lavoura-pecuária? $\circ$ Sim $\circ$ Não --- Produtos/serviços: \_\_\_\_\_\_\_\_ \\
$\Box$ \textbf{Floresta/Mata} & $\Box$ Floresta (semi)natural; $\Box$ Plantação florestal/reflorestamento & Tipo de árvore: \_\_\_\_\_\_\_\_ $\circ$ Decídua $\circ$ Mista $\circ$ Perenifólia --- Produtos: $\Box$ Madeira; $\Box$ Lenha; $\Box$ Frutas/nozes; $\Box$ Produtos não madeireiros; $\Box$ Pastejo; $\Box$ Conservação; $\Box$ Recreação; $\Box$ Proteção contra desastres \\
$\Box$ \textbf{Assentamentos/infraestrutura} & $\Box$ Edificações; $\Box$ Transporte (estradas, ferrovias); $\Box$ Energia (dutos, linhas de transmissão); $\Box$ Outro & Observações: \_\_\_\_\_\_\_\_ \\
$\Box$ \textbf{Cursos d'água/corpos hídricos/áreas úmidas} & $\Box$ Drenagens/canais; $\Box$ Lagoas/barragens; $\Box$ Pântanos/áreas úmidas; $\Box$ Rios e zona ripária; $\Box$ Lagos e margens; $\Box$ Mar e litoral & Produtos/serviços: \_\_\_\_\_\_\_\_ \\
$\Box$ \textbf{Mineração/indústrias extrativas} & Especifique: \_\_\_\_\_\_\_\_ & Produtos: \_\_\_\_\_\_\_\_ \\
$\Box$ \textbf{Terra improdutiva} & Especifique: \_\_\_\_\_\_\_\_ & Observações: \_\_\_\_\_\_\_\_ \\
$\Box$ \textbf{Áreas protegidas} & Especifique: \_\_\_\_\_\_\_\_ & Observações: \_\_\_\_\_\_\_\_ \\
\end{longtable}
}

\subsubsection{3.3 Uso da terra antes da implementação da Tecnologia {[Campo]}}\label{uso-da-terra-antes-da-implementauxe7uxe3o-da-tecnologia}

O uso da terra mudou devido à implementação da Tecnologia?

$\circ$ Não (pule para 3.4) $\circ$ Sim (preencha abaixo com relação ao uso anterior à implementação)

\emph{Se sim, repita as mesmas categorias da tabela acima (3.2) referentes ao uso anterior.}

\subsubsection{3.4 Suprimento hídrico}\label{suprimento-huxeddrico}

Suprimento hídrico para a terra onde a Tecnologia é aplicada:

$\circ$ Sequeiro $\circ$ Misto (sequeiro--irrigado) $\circ$ Irrigação plena $\circ$ Outro (p.ex. pós-inundação): \_\_\_\_\_\_\_\_\_\_\_\_\_\_\_\_\_\_\_\_\_\_\_\_\_\_\_\_\_\_\_\_\_\_\_\_\_\_\_\_

Comentário: \_\_\_\_\_\_\_\_\_\_\_\_\_\_\_\_\_\_\_\_\_\_\_\_\_\_\_\_\_\_\_\_\_\_\_\_\_\_\_\_

\textbf{Definições:}

\begin{itemize}
\tightlist
\item
  \textbf{Sequeiro:} estabelecimento e desenvolvimento da cultura completamente determinados pela precipitação.
\item
  \textbf{Misto sequeiro--irrigado:} aplicação limitada de água quando a precipitação é insuficiente, para aumentar e estabilizar a produtividade.
\item
  \textbf{Irrigação plena:} fornecimento artificial regular de água, além da chuva.
\item
  \textbf{Pós-inundação:} após inundação natural do campo, a água infiltrada é utilizada como reserva hídrica para cultivo.
\end{itemize}

\subsubsection{3.5 Grupo de MST ao qual a Tecnologia pertence}\label{grupo-de-mst-ao-qual-a-tecnologia-pertence}

\emph{Atribua a Tecnologia a um dos seguintes grupos. Se não for possível, selecione no máximo 3 grupos.}

\begin{itemize}
\tightlist
\item[$\square$]
  Manejo de florestas naturais/seminaturais
\item[$\square$]
  Manejo de plantações florestais
\item[$\square$]
  Agrofloresta
\item[$\square$]
  Quebra-vento/cinturão verde
\item[$\square$]
  Fechamento de área (cessar uso, apoiar restauração)
\item[$\square$]
  Sistema rotacional (rotação de culturas, pousio, agricultura itinerante)
\item[$\square$]
  Pastoralismo e manejo de pastagens
\item[$\square$]
  Integração lavoura-pecuária
\item[$\square$]
  Melhoria da cobertura do solo/vegetação
\item[$\square$]
  Distúrbio mínimo do solo
\item[$\square$]
  Manejo integrado da fertilidade do solo
\item[$\square$]
  Medida em nível (\emph{cross-slope})
\item[$\square$]
  Manejo integrado de pragas e doenças (incl.~agricultura orgânica)
\item[$\square$]
  Variedades vegetais/raças animais melhoradas
\item[$\square$]
  Captação de água
\item[$\square$]
  Manejo de irrigação (incl.~abastecimento, drenagem)
\item[$\square$]
  Desvio e drenagem de água
\item[$\square$]
  Manejo de águas superficiais (nascentes, rios, lagos, mar, zona ripária)
\item[$\square$]
  Manejo de águas subterrâneas
\item[$\square$]
  Proteção/manejo de áreas úmidas
\item[$\square$]
  Gestão de resíduos/efluentes
\item[$\square$]
  Eficiência energética
\item[$\square$]
  Apicultura, aquicultura, avicultura, cunicultura, sericicultura, etc.
\item[$\square$]
  Hortas domésticas
\item[$\square$]
  Redução de risco de desastres baseada em ecossistemas
\item[$\square$]
  Medidas pós-colheita
\item[$\square$]
  Outro (especifique): \_\_\_\_\_\_\_\_\_\_\_\_\_\_\_\_\_\_\_\_\_\_\_\_\_\_\_\_\_\_\_\_\_\_\_\_\_\_\_\_
\end{itemize}

\subsubsection{3.6 Medidas de MST que compõem a Tecnologia}\label{medidas-de-mst-que-compuxf5em-a-tecnologia}

\emph{Várias respostas possíveis.}

{
\begin{longtable}[]{@{}
  >{\raggedright\arraybackslash}p{(\linewidth - 4\tabcolsep) * \real{0.3902}}
  >{\raggedright\arraybackslash}p{(\linewidth - 4\tabcolsep) * \real{0.3659}}
  >{\raggedright\arraybackslash}p{(\linewidth - 4\tabcolsep) * \real{0.2439}}@{}}
\toprule\noalign{}
\begin{minipage}[b]{\linewidth}\raggedright
Tipo de medida
\end{minipage} & \begin{minipage}[b]{\linewidth}\raggedright
Subcategorias
\end{minipage} & \begin{minipage}[b]{\linewidth}\raggedright
Exemplos
\end{minipage} \\
\midrule\noalign{}
\endhead
\bottomrule\noalign{}
\endlastfoot
\textbf{Agronômicas} (associadas a cultivos anuais; repetidas sazonalmente; curta duração; sem alteração do perfil do terreno) & A1: Cobertura vegetal/solo; A2: Matéria orgânica/fertilidade; A3: Tratamento superficial do solo; A4: Tratamento subsuperficial; A5: Manejo de sementes, variedades melhoradas; A6: Manejo de resíduos; A7: Outras & Consórcio, cobertura morta, plantio direto, cultivo em nível, compostagem, adubação verde, rotação de culturas \\
\textbf{Vegetativas} (uso de gramíneas perenes, arbustos ou árvores; longa duração) & V1: Cobertura arbórea/arbustiva; V2: Gramíneas e herbáceas perenes; V3: Remoção de vegetação; V4: Substituição/remoção de espécies invasoras; V5: Outras & Agrofloresta, quebra-ventos, faixas de gramíneas em nível, cercas vivas, viveiros \\
\textbf{Estruturais} (permanentes; requerem investimento substancial; movimentação de terra e/ou construção) & S1: Terraços; S2: Cordões/camalhões; S3: Valas graduadas/canais; S4: Valas/covas niveladas; S5: Barragens/tanques; S6: Muros/barreiras/paliçadas/cercas; S7: Captação/irrigação; S8: Saneamento/efluentes; S9: Abrigos para plantas/animais; S10: Eficiência energética; S11: Outras & Terraços de bancada, cordões de pedra, valas de infiltração, barraginhas, estabilização de ravinas (\emph{check dams}), captação de telhado \\
\textbf{De manejo} (mudança fundamental no uso da terra; intensidade reduzida) & M1: Mudança de tipo de uso; M2: Mudança de manejo/intensidade; M3: Disposição conforme ambiente; M4: Mudança de cronograma; M5: Controle de composição de espécies; M6: Gestão de resíduos; M7: Outras & Fechamento de área, pousio gerenciado, acesso controlado, ajuste de carga animal, queima prescrita \\
\textbf{Outras} & --- & Apicultura, aquicultura, armazenamento de alimentos, processamento pós-colheita \\
\end{longtable}
}

\textbf{Sistema de preparo do solo (se relevante):}

$\Box$ Plantio direto $\Box$ Preparo reduzido (\textgreater{} 30\% cobertura) $\Box$ Preparo convencional (\textless{} 30\% cobertura)

\textbf{Manejo de resíduos (se relevante):}

$\Box$ Queimados $\Box$ Pastejados $\Box$ Recolhidos $\Box$ Mantidos

\subsubsection{3.7 Principais tipos de degradação da terra abordados pela Tecnologia {[Campo]}}\label{principais-tipos-de-degradauxe7uxe3o-da-terra-abordados-pela-tecnologia}

\emph{Várias respostas possíveis.}

\subsubsection{W --- Erosão hídrica do solo}\label{w-erosuxe3o-huxeddrica-do-solo}

\begin{itemize}
\tightlist
\item[$\square$]
  Wt: Perda de solo superficial/erosão laminar (remoção uniforme da camada superficial)
\item[$\square$]
  Wg: Erosão em ravinas/voçorocas (remoção ao longo de linhas de drenagem, canais \textgreater{} 30 cm)
\item[$\square$]
  Wm: Movimentos de massa/deslizamentos (queda ou deslizamento de massa de terra em encosta)
\item[$\square$]
  Wr: Erosão de margens fluviais (desgaste das margens de cursos d'água)
\item[$\square$]
  Wc: Erosão costeira (perda de terra ao longo da costa por ação de ondas/correntes/marés)
\item[$\square$]
  Wo: Efeitos fora do local (deposição de sedimentos, enchentes a jusante, assoreamento)
\end{itemize}

\subsubsection{E --- Erosão eólica do solo}\label{e-erosuxe3o-euxf3lica-do-solo}

\begin{itemize}
\tightlist
\item[$\square$]
  Et: Perda superficial (deslocamento uniforme)
\item[$\square$]
  Ed: Deflação e deposição (remoção desigual de material)
\item[$\square$]
  Eo: Efeitos fora do local (cobertura do terreno por partículas transportadas pelo vento)
\end{itemize}

\subsubsection{C --- Deterioração química do solo}\label{c-deteriorauxe7uxe3o-quuxedmica-do-solo}

\begin{itemize}
\tightlist
\item[$\square$]
  Cn: Declínio de fertilidade e redução de matéria orgânica (lixiviação, exaustão)
\item[$\square$]
  Ca: Acidificação (redução do pH do solo)
\item[$\square$]
  Cp: Poluição do solo (contaminação por materiais tóxicos)
\item[$\square$]
  Cs: Salinização/alcalinização (aumento líquido de sais na camada superficial)
\end{itemize}

\subsubsection{P --- Deterioração física do solo}\label{p-deteriorauxe7uxe3o-fuxedsica-do-solo}

\begin{itemize}
\tightlist
\item[$\square$]
  Pc: Compactação (deterioração da estrutura por pisoteio ou maquinário)
\item[$\square$]
  Pk: Selamento superficial/\emph{crusting} (entupimento de poros, camada impermeável superficial)
\item[$\square$]
  Pi: Impermeabilização do solo (cobertura por material impermeável: construção, mineração, estradas)
\item[$\square$]
  Pw: Encharcamento (saturação hídrica induzida por ação humana)
\item[$\square$]
  Ps: Subsidência de solos orgânicos (rebaixamento da superfície)
\item[$\square$]
  Pu: Perda de função bioprodutiva por outras atividades
\end{itemize}

\subsubsection{B --- Degradação biológica}\label{b-degradauxe7uxe3o-bioluxf3gica}

\begin{itemize}
\tightlist
\item[$\square$]
  Bc: Redução de cobertura vegetal (aumento de solo exposto)
\item[$\square$]
  Bh: Perda de habitats (redução de diversidade de vegetação, fragmentação)
\item[$\square$]
  Bq: Declínio de biomassa (redução de produção vegetal)
\item[$\square$]
  Bf: Efeitos deletérios de incêndios
\item[$\square$]
  Bs: Declínio de qualidade/composição/diversidade de espécies
\item[$\square$]
  Bl: Perda de vida do solo (declínio de macro e microrganismos)
\item[$\square$]
  Bp: Aumento de pragas/doenças, perda de predadores
\end{itemize}

\subsubsection{H --- Degradação hídrica}\label{h-degradauxe7uxe3o-huxeddrica}

\begin{itemize}
\tightlist
\item[$\square$]
  Ha: Aridificação (redução da umidade média do solo)
\item[$\square$]
  Hs: Mudança na quantidade de águas superficiais (regime de vazão)
\item[$\square$]
  Hg: Mudança no nível do lençol freático/aquífero
\item[$\square$]
  Hp: Declínio na qualidade de águas superficiais
\item[$\square$]
  Hq: Declínio na qualidade de águas subterrâneas
\item[$\square$]
  Hw: Redução da capacidade tamponante de áreas úmidas
\end{itemize}

\subsubsection{3.8 Prevenção, redução ou restauração da degradação da terra}\label{prevenuxe7uxe3o-reduuxe7uxe3o-ou-restaurauxe7uxe3o-da-degradauxe7uxe3o-da-terra}

\emph{Marque no máximo duas opções.}

\begin{itemize}
\tightlist
\item
  $\circ$ Prevenir/evitar a degradação da terra
\item
  $\circ$ Reduzir a degradação da terra
\item
  $\circ$ Restaurar/reabilitar terra severamente degradada / reverter a degradação
\item
  $\circ$ Adaptar-se à degradação da terra
\item
  $\circ$ Não aplicável
\end{itemize}

Comentários: \_\_\_\_\_\_\_\_\_\_\_\_\_\_\_\_\_\_\_\_\_\_\_\_\_\_\_\_\_\_\_\_\_\_\_\_\_\_\_\_

\textbf{Definições:}

\begin{itemize}
\tightlist
\item
  \textbf{Prevenir (evitar):} uso de boas práticas de manejo em terras que podem ser propensas à degradação, mantendo recursos naturais e suas funções.
\item
  \textbf{Reduzir:} intervenções para reduzir degradação em curso e/ou deter degradação adicional. Impactos perceptíveis no curto a médio prazo.
\item
  \textbf{Reabilitar/restaurar/reverter:} necessário quando a terra está tão degradada que o uso original não é mais possível. Investimentos de longo prazo.
\item
  \textbf{Adaptar:} quando restauração não é mais possível ou está além dos recursos dos usuários. O estado de degradação é ``aceito'' e o manejo é adaptado.
\end{itemize}

\newpage

\begin{center}\rule{0.5\linewidth}{0.5pt}\end{center}

\subsection{4. Especificações técnicas, atividades de implantação, insumos e custos}\label{especificauxe7uxf5es-tuxe9cnicas-atividades-de-implantauxe7uxe3o-insumos-e-custos}

\subsubsection{4.1 Desenho técnico da Tecnologia {[Campo]}}\label{desenho-tuxe9cnico-da-tecnologia}

\emph{Forneça desenho(s) detalhado(s) e abrangente(s) (incluindo dimensões) da Tecnologia e indique especificações técnicas, medições, espaçamento, gradiente, etc. Mantenha o desenho simples e esquemático. O desenho técnico é crucial para a compreensão da Tecnologia.}

\begin{itemize}
\tightlist
\item
  Formatos aceitos: PDF, JPG, PNG (máx. 3 MB)
\item
  Formato quadrado é ideal
\item
  Apenas símbolos e/ou números no desenho; texto explicativo em campo separado
\end{itemize}

Autor: \_\_\_\_\_\_\_\_ Data: \_\_\_\_\_\_\_\_

\textbf{Resumo das especificações técnicas} (dimensões, espaçamento, intervalos, gradiente, capacidade, materiais, espécies, densidade, etc.):

\begin{center}\rule{0.5\linewidth}{0.5pt}\end{center}

\begin{center}\rule{0.5\linewidth}{0.5pt}\end{center}

\begin{center}\rule{0.5\linewidth}{0.5pt}\end{center}

\subsubsection{4.2 Informações gerais para cálculo de insumos e custos}\label{informauxe7uxf5es-gerais-para-cuxe1lculo-de-insumos-e-custos}

\emph{Distingua entre implantação/investimento inicial e manutenção/atividades recorrentes anuais. Todos os custos devem ser calculados com base em preços de mercado. Se a mão de obra é fornecida pelos próprios usuários, indique o custo equivalente de mão de obra contratada.}

Especifique como custos e insumos foram calculados:

$\circ$ Por área da Tecnologia $\rightarrow$ tamanho e unidade de área: \_\_\_\_\_\_\_\_ Se usar unidade local, fator de conversão: 1 hectare = \_\_\_\_\_\_\_\_

$\circ$ Por unidade da Tecnologia $\rightarrow$ especifique unidade: \_\_\_\_\_\_\_\_ Dimensões da unidade: \_\_\_\_\_\_\_\_

Moeda utilizada: $\circ$ Dólares americanos (USD) $\circ$ Outra (especifique): \_\_\_\_\_\_\_\_

Taxa de câmbio (se relevante): 1 USD = \_\_\_\_\_\_\_\_

Custo médio diário de mão de obra contratada: \_\_\_\_\_\_\_\_

\subsubsection{4.3 Atividades de implantação {[Campo]}}\label{atividades-de-implantauxe7uxe3o}

\emph{Liste as atividades de implantação na sequência e indique o período.}

{
\begin{longtable}[]{@{}lll@{}}
\toprule\noalign{}
Nº & Atividade & Período \\
\midrule\noalign{}
\endhead
\bottomrule\noalign{}
\endlastfoot
1 & \_\_\_\_\_\_\_\_\_\_\_\_\_\_\_\_\_\_\_\_\_\_\_\_\_\_\_\_\_\_\_\_\_\_\_\_\_\_\_\_ & \_\_\_\_\_\_\_\_ \\
2 & \_\_\_\_\_\_\_\_\_\_\_\_\_\_\_\_\_\_\_\_\_\_\_\_\_\_\_\_\_\_\_\_\_\_\_\_\_\_\_\_ & \_\_\_\_\_\_\_\_ \\
3 & \_\_\_\_\_\_\_\_\_\_\_\_\_\_\_\_\_\_\_\_\_\_\_\_\_\_\_\_\_\_\_\_\_\_\_\_\_\_\_\_ & \_\_\_\_\_\_\_\_ \\
4 & \_\_\_\_\_\_\_\_\_\_\_\_\_\_\_\_\_\_\_\_\_\_\_\_\_\_\_\_\_\_\_\_\_\_\_\_\_\_\_\_ & \_\_\_\_\_\_\_\_ \\
5 & \_\_\_\_\_\_\_\_\_\_\_\_\_\_\_\_\_\_\_\_\_\_\_\_\_\_\_\_\_\_\_\_\_\_\_\_\_\_\_\_ & \_\_\_\_\_\_\_\_ \\
6 & \_\_\_\_\_\_\_\_\_\_\_\_\_\_\_\_\_\_\_\_\_\_\_\_\_\_\_\_\_\_\_\_\_\_\_\_\_\_\_\_ & \_\_\_\_\_\_\_\_ \\
\end{longtable}
}

\textbf{Período:} época em que a atividade é realizada, p.ex. mês ou estação, ou ``após a colheita'', ``antes do início das chuvas'', etc.

\subsubsection{4.4 Custos dos insumos para implantação {[Campo]}}\label{custos-dos-insumos-para-implantauxe7uxe3o}

\emph{Custos devem referir-se à área/unidade da Tecnologia definida em 4.2 e às atividades listadas em 4.3.}

{
\begin{longtable}[]{@{}
  >{\raggedright\arraybackslash}p{(\linewidth - 12\tabcolsep) * \real{0.0952}}
  >{\raggedright\arraybackslash}p{(\linewidth - 12\tabcolsep) * \real{0.1786}}
  >{\raggedright\arraybackslash}p{(\linewidth - 12\tabcolsep) * \real{0.1071}}
  >{\raggedleft\arraybackslash}p{(\linewidth - 12\tabcolsep) * \real{0.0833}}
  >{\raggedleft\arraybackslash}p{(\linewidth - 12\tabcolsep) * \real{0.1310}}
  >{\raggedleft\arraybackslash}p{(\linewidth - 12\tabcolsep) * \real{0.1548}}
  >{\centering\arraybackslash}p{(\linewidth - 12\tabcolsep) * \real{0.2500}}@{}}
\toprule\noalign{}
\begin{minipage}[b]{\linewidth}\raggedright
Insumo
\end{minipage} & \begin{minipage}[b]{\linewidth}\raggedright
Especificação
\end{minipage} & \begin{minipage}[b]{\linewidth}\raggedright
Unidade
\end{minipage} & \begin{minipage}[b]{\linewidth}\raggedleft
Qtde.
\end{minipage} & \begin{minipage}[b]{\linewidth}\raggedleft
Custo/un.
\end{minipage} & \begin{minipage}[b]{\linewidth}\raggedleft
Custo total
\end{minipage} & \begin{minipage}[b]{\linewidth}\centering
\% arc. pelo usuário
\end{minipage} \\
\midrule\noalign{}
\endhead
\bottomrule\noalign{}
\endlastfoot
\textbf{Mão de obra} & & & & & & \\
& & & & & & \\
\textbf{Equipamento} & & & & & & \\
& & & & & & \\
\textbf{Material vegetal} & & & & & & \\
& & & & & & \\
\textbf{Fertilizantes e biocidas} & & & & & & \\
& & & & & & \\
\textbf{Material de construção} & & & & & & \\
& & & & & & \\
\textbf{Outros} & & & & & & \\
& & & & & & \\
\textbf{CUSTO TOTAL DE IMPLANTAÇÃO} & & & & & \textbf{\_\_\_\_\_\_\_\_} & \\
\end{longtable}
}

\textbf{Notas sobre insumos:}

\begin{itemize}
\tightlist
\item
  \textbf{Mão de obra:} total de pessoas-dia (pagas ou não). Sob ``Custo/un.'', indicar diária.
\item
  \textbf{Equipamento:} ferramentas, horas-máquina, tração animal. Calcular com base em custos de aluguel.
\item
  \textbf{Material vegetal:} sementes, mudas, estacas, etc.
\item
  \textbf{Fertilizantes e biocidas:} composto/esterco, fertilizante inorgânico, herbicidas, pesticidas, etc.
\item
  \textbf{Material de construção:} madeira, pedras, terra, cimento, tubos, tanques, etc.
\item
  \textbf{\% arcado pelo usuário:} percentual de custos que os usuários contribuem.
\end{itemize}

Se não for possível detalhar, estime o custo total de implantação: \_\_\_\_\_\_\_\_

\subsubsection{4.5 Atividades de manutenção/recorrentes {[Campo]}}\label{atividades-de-manutenuxe7uxe3orecorrentes}

{
\begin{longtable}[]{@{}lll@{}}
\toprule\noalign{}
Nº & Atividade & Período/Frequência \\
\midrule\noalign{}
\endhead
\bottomrule\noalign{}
\endlastfoot
1 & \_\_\_\_\_\_\_\_\_\_\_\_\_\_\_\_\_\_\_\_\_\_\_\_\_\_\_\_\_\_\_\_\_\_\_\_\_\_\_\_ & \_\_\_\_\_\_\_\_ \\
2 & \_\_\_\_\_\_\_\_\_\_\_\_\_\_\_\_\_\_\_\_\_\_\_\_\_\_\_\_\_\_\_\_\_\_\_\_\_\_\_\_ & \_\_\_\_\_\_\_\_ \\
3 & \_\_\_\_\_\_\_\_\_\_\_\_\_\_\_\_\_\_\_\_\_\_\_\_\_\_\_\_\_\_\_\_\_\_\_\_\_\_\_\_ & \_\_\_\_\_\_\_\_ \\
4 & \_\_\_\_\_\_\_\_\_\_\_\_\_\_\_\_\_\_\_\_\_\_\_\_\_\_\_\_\_\_\_\_\_\_\_\_\_\_\_\_ & \_\_\_\_\_\_\_\_ \\
5 & \_\_\_\_\_\_\_\_\_\_\_\_\_\_\_\_\_\_\_\_\_\_\_\_\_\_\_\_\_\_\_\_\_\_\_\_\_\_\_\_ & \_\_\_\_\_\_\_\_ \\
\end{longtable}
}

\subsubsection{4.6 Custos dos insumos para manutenção (por ano) {[Campo]}}\label{custos-dos-insumos-para-manutenuxe7uxe3o-por-ano}

{
\begin{longtable}[]{@{}
  >{\raggedright\arraybackslash}p{(\linewidth - 12\tabcolsep) * \real{0.0952}}
  >{\raggedright\arraybackslash}p{(\linewidth - 12\tabcolsep) * \real{0.1786}}
  >{\raggedright\arraybackslash}p{(\linewidth - 12\tabcolsep) * \real{0.1071}}
  >{\raggedleft\arraybackslash}p{(\linewidth - 12\tabcolsep) * \real{0.0833}}
  >{\raggedleft\arraybackslash}p{(\linewidth - 12\tabcolsep) * \real{0.1310}}
  >{\raggedleft\arraybackslash}p{(\linewidth - 12\tabcolsep) * \real{0.1548}}
  >{\centering\arraybackslash}p{(\linewidth - 12\tabcolsep) * \real{0.2500}}@{}}
\toprule\noalign{}
\begin{minipage}[b]{\linewidth}\raggedright
Insumo
\end{minipage} & \begin{minipage}[b]{\linewidth}\raggedright
Especificação
\end{minipage} & \begin{minipage}[b]{\linewidth}\raggedright
Unidade
\end{minipage} & \begin{minipage}[b]{\linewidth}\raggedleft
Qtde.
\end{minipage} & \begin{minipage}[b]{\linewidth}\raggedleft
Custo/un.
\end{minipage} & \begin{minipage}[b]{\linewidth}\raggedleft
Custo total
\end{minipage} & \begin{minipage}[b]{\linewidth}\centering
\% arc. pelo usuário
\end{minipage} \\
\midrule\noalign{}
\endhead
\bottomrule\noalign{}
\endlastfoot
\textbf{Mão de obra} & & & & & & \\
\textbf{Equipamento} & & & & & & \\
\textbf{Material vegetal} & & & & & & \\
\textbf{Fertilizantes e biocidas} & & & & & & \\
\textbf{Material de construção} & & & & & & \\
\textbf{Outros} & & & & & & \\
\textbf{CUSTO TOTAL DE MANUTENÇÃO} & & & & & \textbf{\_\_\_\_\_\_\_\_} & \\
\end{longtable}
}

\subsubsection{4.7 Fatores mais importantes que afetam os custos {[Campo]}}\label{fatores-mais-importantes-que-afetam-os-custos}

\begin{center}\rule{0.5\linewidth}{0.5pt}\end{center}

\begin{center}\rule{0.5\linewidth}{0.5pt}\end{center}

\newpage

\begin{center}\rule{0.5\linewidth}{0.5pt}\end{center}

\subsection{5. Ambiente natural e humano}\label{ambiente-natural-e-humano}

\emph{Forneça detalhes das condições naturais (biofísicas) onde a Tecnologia é aplicada. Descreva as condições tal como eram sem qualquer impacto do manejo sustentável da terra.}

\subsubsection{5.1 Clima}\label{clima}

\emph{Marque no máximo duas respostas por pergunta.}

\textbf{Precipitação anual:}

$\circ$ \textless{} 250 mm $\circ$ 251--500 mm $\circ$ 501--750 mm $\circ$ 751--1.000 mm $\circ$ 1.001--1.500 mm $\circ$ 1.501--2.000 mm $\circ$ 2.001--3.000 mm $\circ$ 3.001--4.000 mm $\circ$ \textgreater{} 4.000 mm

Precipitação média anual (se conhecida): \_\_\_\_\_\_\_\_ mm

\textbf{Zona agroclimática:}

$\circ$ Úmida (período de crescimento \textgreater{} 270 dias) $\circ$ Subúmida (180--269 dias) $\circ$ Semiárida (75--179 dias) $\circ$ Árida (\textless{} 74 dias)

\emph{Período de crescimento (LGP): período em que a precipitação excede metade da evapotranspiração potencial (ETP) e a temperatura é superior a 6,5°C.}

Especificações/comentários sobre distribuição da precipitação, sazonalidade, estações chuvosas, chuvas intensas, períodos secos: \_\_\_\_\_\_\_\_\_\_\_\_\_\_\_\_\_\_\_\_\_\_\_\_\_\_\_\_\_\_\_\_\_\_\_\_\_\_\_\_

Estação meteorológica de referência: \_\_\_\_\_\_\_\_\_\_\_\_\_\_\_\_\_\_\_\_\_\_\_\_\_\_\_\_\_\_\_\_\_\_\_\_\_\_\_\_

Especificações sobre clima (p.ex. temperatura média anual): \_\_\_\_\_\_\_\_\_\_\_\_\_\_\_\_\_\_\_\_\_\_\_\_\_\_\_\_\_\_\_\_\_\_\_\_\_\_\_\_

\subsubsection{5.2 Topografia {[Campo]}}\label{topografia}

\emph{Marque no máximo duas respostas por pergunta.}

\textbf{Declividade média:}

$\circ$ Plana (0--2\%) $\circ$ Suave (3--5\%) $\circ$ Moderada (6--10\%) $\circ$ Ondulada (11--15\%) $\circ$ Colinosa (16--30\%) $\circ$ Íngreme (31--60\%) $\circ$ Muito íngreme (\textgreater{} 60\%)

\textbf{Formas de relevo:}

$\circ$ Planalto/planícies $\circ$ Cumeeiras $\circ$ Encostas de montanha $\circ$ Encostas de colina $\circ$ Sopé de encosta $\circ$ Fundos de vale

\textbf{Zona altitudinal:}

$\circ$ \textless{} 100 m $\circ$ 101--500 m $\circ$ 501--1.000 m $\circ$ 1.001--1.500 m $\circ$ 1.501--2.000 m $\circ$ 2.001--2.500 m $\circ$ 2.501--3.000 m $\circ$ 3.001--4.000 m $\circ$ \textgreater{} 4.000 m

A Tecnologia é especificamente aplicada em:

$\circ$ Situações convexas (cumeeira --- divergência do fluxo de água) $\circ$ Situações côncavas (depressão --- convergência do fluxo) $\circ$ Não relevante

Comentários sobre topografia: \_\_\_\_\_\_\_\_\_\_\_\_\_\_\_\_\_\_\_\_\_\_\_\_\_\_\_\_\_\_\_\_\_\_\_\_\_\_\_\_

\subsubsection{5.3 Solos {[Campo]}}\label{solos}

\emph{Parâmetros baseados nos padrões FAO. Marque no máximo duas respostas por pergunta.}

\textbf{Profundidade média do solo:}

$\circ$ Muito raso (0--20 cm) $\circ$ Raso (21--50 cm) $\circ$ Moderadamente profundo (51--80 cm) $\circ$ Profundo (81--120 cm) $\circ$ Muito profundo (\textgreater{} 120 cm)

\textbf{Textura do solo (horizonte superficial):}

$\circ$ Grossa/leve (arenosa) $\circ$ Média (franca, siltosa) $\circ$ Fina/pesada (argilosa)

\textbf{Textura do solo (\textgreater{} 20 cm abaixo da superfície):}

$\circ$ Grossa/leve (arenosa) $\circ$ Média (franca, siltosa) $\circ$ Fina/pesada (argilosa)

\textbf{Matéria orgânica no horizonte superficial:}

$\circ$ Alta (\textgreater{} 3\%) $\circ$ Média (1--3\%) $\circ$ Baixa (\textless{} 1\%)

Informações adicionais sobre o solo (tipo, pH, CTC, nitrogênio, salinidade, etc.): \_\_\_\_\_\_\_\_\_\_\_\_\_\_\_\_\_\_\_\_\_\_\_\_\_\_\_\_\_\_\_\_\_\_\_\_\_\_\_\_

\subsubsection{5.4 Disponibilidade e qualidade da água {[Campo]}}\label{disponibilidade-e-qualidade-da-uxe1gua}

\emph{Uma resposta por pergunta.}

\textbf{Nível do lençol freático:}

$\circ$ Superficial $\circ$ \textless{} 5 m $\circ$ 5--50 m $\circ$ \textgreater{} 50 m

\textbf{Disponibilidade de águas superficiais:}

$\circ$ Excesso (p.ex. encharcamento frequente, alto escoamento) $\circ$ Boa (disponível o ano todo) $\circ$ Média (não disponível o ano todo) $\circ$ Pouca/nenhuma

\textbf{Qualidade da água (não tratada):}

$\circ$ Boa (potável) $\circ$ Pobre (requer tratamento) $\circ$ Apenas para uso agrícola (irrigação) $\circ$ Inutilizável

Qualidade refere-se a: $\circ$ Água subterrânea $\circ$ Água superficial $\circ$ Ambas

Salinidade é um problema? $\circ$ Sim $\circ$ Não --- Especifique: \_\_\_\_\_\_\_\_

Inundação da área ocorre? $\circ$ Sim $\circ$ Não --- Se sim: $\circ$ Frequente $\circ$ Episódica

Comentários: \_\_\_\_\_\_\_\_\_\_\_\_\_\_\_\_\_\_\_\_\_\_\_\_\_\_\_\_\_\_\_\_\_\_\_\_\_\_\_\_

\subsubsection{5.5 Biodiversidade}\label{biodiversidade}

\emph{Uma resposta por pergunta.}

\textbf{Diversidade de espécies:} $\circ$ Alta $\circ$ Média $\circ$ Baixa

\textbf{Diversidade de habitats:} $\circ$ Alta $\circ$ Média $\circ$ Baixa

Comentários: \_\_\_\_\_\_\_\_\_\_\_\_\_\_\_\_\_\_\_\_\_\_\_\_\_\_\_\_\_\_\_\_\_\_\_\_\_\_\_\_

\subsubsection{5.6 Características dos usuários da terra que aplicam a Tecnologia}\label{caracteruxedsticas-dos-usuuxe1rios-da-terra-que-aplicam-a-tecnologia}

\emph{Marque no máximo duas respostas por pergunta.}

\textbf{Sedentário ou nômade:}

$\circ$ Sedentário $\circ$ Seminômade $\circ$ Nômade $\circ$ Outro: \_\_\_\_\_\_\_\_

\textbf{Orientação de mercado do sistema de produção:}

$\circ$ Subsistência (autoconsumo) $\circ$ Misto (subsistência/comercial) $\circ$ Comercial/mercado

\textbf{Nível relativo de riqueza (referência local):}

$\circ$ Muito pobre $\circ$ Pobre $\circ$ Médio $\circ$ Rico $\circ$ Muito rico

\textbf{Individual ou grupo:}

$\circ$ Individual/domicílio $\circ$ Grupos/comunidade $\circ$ Cooperativa $\circ$ Empregado (empresa, governo)

\textbf{Gênero:} $\Box$ Mulheres $\Box$ Homens

\textbf{Faixa etária:} $\Box$ Crianças $\Box$ Jovens $\Box$ Meia-idade $\Box$ Idosos

\textbf{Renda extra-propriedade:}

$\circ$ Menos de 10\% da renda total $\circ$ 10--50\% $\circ$ \textgreater{} 50\%

\textbf{Nível de mecanização:}

$\circ$ Trabalho manual $\circ$ Tração animal $\circ$ Mecanizado/motorizado

Outras características relevantes: \_\_\_\_\_\_\_\_\_\_\_\_\_\_\_\_\_\_\_\_\_\_\_\_\_\_\_\_\_\_\_\_\_\_\_\_\_\_\_\_

\subsubsection{5.7 Área média possuída, arrendada ou utilizada (com direitos de uso) pelos usuários {[Campo]}}\label{uxe1rea-muxe9dia-possuuxedda-arrendada-ou-utilizada-com-direitos-de-uso-pelos-usuuxe1rios}

\emph{Marque no máximo duas respostas.}

$\circ$ \textless{} 0,5 ha $\circ$ 0,5--1 ha $\circ$ 1--2 ha $\circ$ 2--5 ha $\circ$ 5--15 ha $\circ$ 15--50 ha $\circ$ 50--100 ha $\circ$ 100--500 ha $\circ$ 500--1.000 ha $\circ$ 1.000--10.000 ha $\circ$ \textgreater{} 10.000 ha

Considerada (no contexto local): $\circ$ Pequena escala $\circ$ Média escala $\circ$ Grande escala

\subsubsection{5.8 Propriedade da terra, direitos de uso da terra e direitos de uso da água {[Campo]}}\label{propriedade-da-terra-direitos-de-uso-da-terra-e-direitos-de-uso-da-uxe1gua}

\emph{Marque no máximo duas respostas por pergunta.}

\textbf{Propriedade da terra:}

$\circ$ Estado $\circ$ Empresa $\circ$ Comunal/vilarejo $\circ$ Grupo $\circ$ Individual sem título $\circ$ Individual com título $\circ$ Outro: \_\_\_\_\_\_\_\_

\textbf{Direitos de uso da terra:}

$\circ$ Acesso livre (não organizado) $\circ$ Comunal (organizado) $\circ$ Arrendado $\circ$ Individual $\circ$ Outro: \_\_\_\_\_\_\_\_

\textbf{Direitos de uso da água (se relevante):}

$\circ$ Acesso livre (não organizado) $\circ$ Comunal (organizado) $\circ$ Arrendado $\circ$ Individual $\circ$ Outro: \_\_\_\_\_\_\_\_

Direitos de uso baseados em sistema jurídico tradicional?

$\circ$ Sim, especifique: \_\_\_\_\_\_\_\_\_\_\_\_\_\_\_\_\_\_\_\_\_\_\_\_\_\_\_\_\_\_\_\_\_\_\_\_\_\_\_\_ $\circ$ Não, especifique: \_\_\_\_\_\_\_\_\_\_\_\_\_\_\_\_\_\_\_\_\_\_\_\_\_\_\_\_\_\_\_\_\_\_\_\_\_\_\_\_

\subsubsection{5.9 Acesso a serviços e infraestrutura}\label{acesso-a-serviuxe7os-e-infraestrutura}

\emph{Várias respostas possíveis.}

{
\begin{longtable}[]{@{}lccc@{}}
\toprule\noalign{}
Serviço & Precário & Moderado & Bom \\
\midrule\noalign{}
\endhead
\bottomrule\noalign{}
\endlastfoot
Saúde & $\circ$ & $\circ$ & $\circ$ \\
Educação & $\circ$ & $\circ$ & $\circ$ \\
Assistência técnica & $\circ$ & $\circ$ & $\circ$ \\
Emprego (p.ex. extra-propriedade) & $\circ$ & $\circ$ & $\circ$ \\
Mercados & $\circ$ & $\circ$ & $\circ$ \\
Energia & $\circ$ & $\circ$ & $\circ$ \\
Estradas e transporte & $\circ$ & $\circ$ & $\circ$ \\
Água potável e saneamento & $\circ$ & $\circ$ & $\circ$ \\
Serviços financeiros & $\circ$ & $\circ$ & $\circ$ \\
Outro: \_\_\_\_\_\_\_\_ & $\circ$ & $\circ$ & $\circ$ \\
\end{longtable}
}

\newpage

\begin{center}\rule{0.5\linewidth}{0.5pt}\end{center}

\subsection{6. Impactos e declarações conclusivas}\label{impactos-e-declarauxe7uxf5es-conclusivas}

\emph{Avalie os impactos relevantes. Se dados baseados em medições não estiverem disponíveis, forneça sua melhor estimativa. ``Negligenciável'' significa ``sem benefício ou desvantagem significativa''. Utilize as colunas ``Quantifique antes/depois do MST'' e ``Comentários'' para demonstrar evidências e justificar sua seleção.}

\subsubsection{\texorpdfstring{6.1 Impactos no local (\emph{on-site}) que a Tecnologia demonstrou}{6.1 Impactos no local (on-site) que a Tecnologia demonstrou}}\label{impactos-no-local-on-site-que-a-tecnologia-demonstrou}

\subsubsection{Impactos socioeconômicos {[Consulta]}}\label{impactos-socioeconuxf4micos}

\textbf{Produção:}

{
\begin{longtable}[]{@{}
  >{\raggedright\arraybackslash}p{(\linewidth - 20\tabcolsep) * \real{0.1286}}
  >{\centering\arraybackslash}p{(\linewidth - 20\tabcolsep) * \real{0.0571}}
  >{\centering\arraybackslash}p{(\linewidth - 20\tabcolsep) * \real{0.0571}}
  >{\centering\arraybackslash}p{(\linewidth - 20\tabcolsep) * \real{0.0571}}
  >{\centering\arraybackslash}p{(\linewidth - 20\tabcolsep) * \real{0.0571}}
  >{\centering\arraybackslash}p{(\linewidth - 20\tabcolsep) * \real{0.0571}}
  >{\centering\arraybackslash}p{(\linewidth - 20\tabcolsep) * \real{0.0571}}
  >{\centering\arraybackslash}p{(\linewidth - 20\tabcolsep) * \real{0.0571}}
  >{\centering\arraybackslash}p{(\linewidth - 20\tabcolsep) * \real{0.1571}}
  >{\centering\arraybackslash}p{(\linewidth - 20\tabcolsep) * \real{0.1429}}
  >{\raggedright\arraybackslash}p{(\linewidth - 20\tabcolsep) * \real{0.1714}}@{}}
\toprule\noalign{}
\begin{minipage}[b]{\linewidth}\raggedright
Impacto
\end{minipage} & \begin{minipage}[b]{\linewidth}\centering
Muito neg.
\end{minipage} & \begin{minipage}[b]{\linewidth}\centering
Negativo
\end{minipage} & \begin{minipage}[b]{\linewidth}\centering
Lev. neg.
\end{minipage} & \begin{minipage}[b]{\linewidth}\centering
Negligenc.
\end{minipage} & \begin{minipage}[b]{\linewidth}\centering
Lev. pos.
\end{minipage} & \begin{minipage}[b]{\linewidth}\centering
Positivo
\end{minipage} & \begin{minipage}[b]{\linewidth}\centering
Muito pos.
\end{minipage} & \begin{minipage}[b]{\linewidth}\centering
Antes MST
\end{minipage} & \begin{minipage}[b]{\linewidth}\centering
Após MST
\end{minipage} & \begin{minipage}[b]{\linewidth}\raggedright
Comentários
\end{minipage} \\
\midrule\noalign{}
\endhead
\bottomrule\noalign{}
\endlastfoot
Produção agrícola (diminuiu $\leftarrow$ $\rightarrow$ aumentou) & $\circ$ & $\circ$ & $\circ$ & $\circ$ & $\circ$ & $\circ$ & $\circ$ & & & \\
Qualidade das culturas & $\circ$ & $\circ$ & $\circ$ & $\circ$ & $\circ$ & $\circ$ & $\circ$ & & & \\
Produção de forragem & $\circ$ & $\circ$ & $\circ$ & $\circ$ & $\circ$ & $\circ$ & $\circ$ & & & \\
Qualidade da forragem & $\circ$ & $\circ$ & $\circ$ & $\circ$ & $\circ$ & $\circ$ & $\circ$ & & & \\
Produção animal & $\circ$ & $\circ$ & $\circ$ & $\circ$ & $\circ$ & $\circ$ & $\circ$ & & & \\
Produção de madeira & $\circ$ & $\circ$ & $\circ$ & $\circ$ & $\circ$ & $\circ$ & $\circ$ & & & \\
Qualidade florestal & $\circ$ & $\circ$ & $\circ$ & $\circ$ & $\circ$ & $\circ$ & $\circ$ & & & \\
Prod. florestal não madeireira & $\circ$ & $\circ$ & $\circ$ & $\circ$ & $\circ$ & $\circ$ & $\circ$ & & & \\
Risco de falha da produção (aum. $\leftarrow$ $\rightarrow$ dim.) & $\circ$ & $\circ$ & $\circ$ & $\circ$ & $\circ$ & $\circ$ & $\circ$ & & & \\
Diversidade de produtos & $\circ$ & $\circ$ & $\circ$ & $\circ$ & $\circ$ & $\circ$ & $\circ$ & & & \\
Área de produção & $\circ$ & $\circ$ & $\circ$ & $\circ$ & $\circ$ & $\circ$ & $\circ$ & & & \\
Manejo da terra (dificultou $\leftarrow$ $\rightarrow$ simplificou) & $\circ$ & $\circ$ & $\circ$ & $\circ$ & $\circ$ & $\circ$ & $\circ$ & & & \\
Geração de energia & $\circ$ & $\circ$ & $\circ$ & $\circ$ & $\circ$ & $\circ$ & $\circ$ & & & \\
\end{longtable}
}

\textbf{Disponibilidade e qualidade da água:}

{
\begin{longtable}[]{@{}lcccl@{}}
\toprule\noalign{}
Impacto & Escala 7 pontos & Antes & Após & Comentários \\
\midrule\noalign{}
\endhead
\bottomrule\noalign{}
\endlastfoot
Disponibilidade de água potável & $\circ$$\circ$$\circ$$\circ$$\circ$$\circ$$\circ$ & & & \\
Qualidade de água potável & $\circ$$\circ$$\circ$$\circ$$\circ$$\circ$$\circ$ & & & \\
Água para pecuária (disponib.) & $\circ$$\circ$$\circ$$\circ$$\circ$$\circ$$\circ$ & & & \\
Água para pecuária (qualidade) & $\circ$$\circ$$\circ$$\circ$$\circ$$\circ$$\circ$ & & & \\
Água para irrigação (disponib.) & $\circ$$\circ$$\circ$$\circ$$\circ$$\circ$$\circ$ & & & \\
Água para irrigação (qualidade) & $\circ$$\circ$$\circ$$\circ$$\circ$$\circ$$\circ$ & & & \\
Demanda de irrigação (aum. $\leftarrow$ $\rightarrow$ dim.) & $\circ$$\circ$$\circ$$\circ$$\circ$$\circ$$\circ$ & & & \\
\end{longtable}
}

\textbf{Renda e custos:}

{
\begin{longtable}[]{@{}lcccl@{}}
\toprule\noalign{}
Impacto & Escala 7 pontos & Antes & Após & Comentários \\
\midrule\noalign{}
\endhead
\bottomrule\noalign{}
\endlastfoot
Despesas com insumos (aum. $\leftarrow$ $\rightarrow$ reduz.) & $\circ$$\circ$$\circ$$\circ$$\circ$$\circ$$\circ$ & & & \\
Renda agrícola & $\circ$$\circ$$\circ$$\circ$$\circ$$\circ$$\circ$ & & & \\
Diversidade de fontes de renda & $\circ$$\circ$$\circ$$\circ$$\circ$$\circ$$\circ$ & & & \\
Disparidades econômicas (aum. $\leftarrow$ $\rightarrow$ dim.) & $\circ$$\circ$$\circ$$\circ$$\circ$$\circ$$\circ$ & & & \\
Carga de trabalho (aum. $\leftarrow$ $\rightarrow$ dim.) & $\circ$$\circ$$\circ$$\circ$$\circ$$\circ$$\circ$ & & & \\
\end{longtable}
}

\subsubsection{Impactos socioculturais {[Consulta]}}\label{impactos-socioculturais}

{
\begin{longtable}[]{@{}lcccl@{}}
\toprule\noalign{}
Impacto & Escala 7 pontos & Antes & Após & Comentários \\
\midrule\noalign{}
\endhead
\bottomrule\noalign{}
\endlastfoot
Segurança alimentar/autossuficiência & $\circ$$\circ$$\circ$$\circ$$\circ$$\circ$$\circ$ & & & \\
Situação de saúde & $\circ$$\circ$$\circ$$\circ$$\circ$$\circ$$\circ$ & & & \\
Direitos de uso da terra/água & $\circ$$\circ$$\circ$$\circ$$\circ$$\circ$$\circ$ & & & \\
Oportunidades culturais & $\circ$$\circ$$\circ$$\circ$$\circ$$\circ$$\circ$ & & & \\
Oportunidades de lazer & $\circ$$\circ$$\circ$$\circ$$\circ$$\circ$$\circ$ & & & \\
Instituições comunitárias & $\circ$$\circ$$\circ$$\circ$$\circ$$\circ$$\circ$ & & & \\
Instituições nacionais & $\circ$$\circ$$\circ$$\circ$$\circ$$\circ$$\circ$ & & & \\
Conhecimento em MST/degradação & $\circ$$\circ$$\circ$$\circ$$\circ$$\circ$$\circ$ & & & \\
Mitigação de conflitos & $\circ$$\circ$$\circ$$\circ$$\circ$$\circ$$\circ$ & & & \\
Situação de grupos desfavorecidos & $\circ$$\circ$$\circ$$\circ$$\circ$$\circ$$\circ$ & & & \\
\end{longtable}
}

\subsubsection{Impactos ecológicos {[Campo]}}\label{impactos-ecoluxf3gicos}

\textbf{Ciclo hidrológico/escoamento:}

{
\begin{longtable}[]{@{}lcccl@{}}
\toprule\noalign{}
Impacto & Escala 7 pontos & Antes & Após & Comentários \\
\midrule\noalign{}
\endhead
\bottomrule\noalign{}
\endlastfoot
Quantidade de água & $\circ$$\circ$$\circ$$\circ$$\circ$$\circ$$\circ$ & & & \\
Qualidade da água & $\circ$$\circ$$\circ$$\circ$$\circ$$\circ$$\circ$ & & & \\
Captação/coleta de água & $\circ$$\circ$$\circ$$\circ$$\circ$$\circ$$\circ$ & & & \\
Escoamento superficial (aum. $\leftarrow$ $\rightarrow$ dim.) & $\circ$$\circ$$\circ$$\circ$$\circ$$\circ$$\circ$ & & & \\
Drenagem & $\circ$$\circ$$\circ$$\circ$$\circ$$\circ$$\circ$ & & & \\
Lençol freático/aquífero & $\circ$$\circ$$\circ$$\circ$$\circ$$\circ$$\circ$ & & & \\
Evaporação (aum. $\leftarrow$ $\rightarrow$ dim.) & $\circ$$\circ$$\circ$$\circ$$\circ$$\circ$$\circ$ & & & \\
\end{longtable}
}

\textbf{Solo:}

{
\begin{longtable}[]{@{}lcccl@{}}
\toprule\noalign{}
Impacto & Escala 7 pontos & Antes & Após & Comentários \\
\midrule\noalign{}
\endhead
\bottomrule\noalign{}
\endlastfoot
Umidade do solo & $\circ$$\circ$$\circ$$\circ$$\circ$$\circ$$\circ$ & & & \\
Cobertura do solo & $\circ$$\circ$$\circ$$\circ$$\circ$$\circ$$\circ$ & & & \\
Perda de solo (aum. $\leftarrow$ $\rightarrow$ dim.) & $\circ$$\circ$$\circ$$\circ$$\circ$$\circ$$\circ$ & & & \\
Acúmulo de solo & $\circ$$\circ$$\circ$$\circ$$\circ$$\circ$$\circ$ & & & \\
Selamento/\emph{crusting} (aum. $\leftarrow$ $\rightarrow$ dim.) & $\circ$$\circ$$\circ$$\circ$$\circ$$\circ$$\circ$ & & & \\
Compactação (aum. $\leftarrow$ $\rightarrow$ dim.) & $\circ$$\circ$$\circ$$\circ$$\circ$$\circ$$\circ$ & & & \\
Ciclagem de nutrientes & $\circ$$\circ$$\circ$$\circ$$\circ$$\circ$$\circ$ & & & \\
Salinidade (aum. $\leftarrow$ $\rightarrow$ dim.) & $\circ$$\circ$$\circ$$\circ$$\circ$$\circ$$\circ$ & & & \\
MOS / carbono subsuperficial & $\circ$$\circ$$\circ$$\circ$$\circ$$\circ$$\circ$ & & & \\
Acidez (aum. $\leftarrow$ $\rightarrow$ dim.) & $\circ$$\circ$$\circ$$\circ$$\circ$$\circ$$\circ$ & & & \\
\end{longtable}
}

\textbf{Biodiversidade (vegetação, animais):}

{
\begin{longtable}[]{@{}lcccl@{}}
\toprule\noalign{}
Impacto & Escala 7 pontos & Antes & Após & Comentários \\
\midrule\noalign{}
\endhead
\bottomrule\noalign{}
\endlastfoot
Cobertura vegetal & $\circ$$\circ$$\circ$$\circ$$\circ$$\circ$$\circ$ & & & \\
Biomassa / C aéreo & $\circ$$\circ$$\circ$$\circ$$\circ$$\circ$$\circ$ & & & \\
Diversidade vegetal & $\circ$$\circ$$\circ$$\circ$$\circ$$\circ$$\circ$ & & & \\
Espécies invasoras (aum. $\leftarrow$ $\rightarrow$ dim.) & $\circ$$\circ$$\circ$$\circ$$\circ$$\circ$$\circ$ & & & \\
Diversidade animal & $\circ$$\circ$$\circ$$\circ$$\circ$$\circ$$\circ$ & & & \\
Espécies benéficas & $\circ$$\circ$$\circ$$\circ$$\circ$$\circ$$\circ$ & & & \\
Espécies nocivas (aum. $\leftarrow$ $\rightarrow$ dim.) & $\circ$$\circ$$\circ$$\circ$$\circ$$\circ$$\circ$ & & & \\
Diversidade de habitats & $\circ$$\circ$$\circ$$\circ$$\circ$$\circ$$\circ$ & & & \\
Pragas/doenças (dim. $\leftarrow$ $\rightarrow$ aum.) & $\circ$$\circ$$\circ$$\circ$$\circ$$\circ$$\circ$ & & & \\
\end{longtable}
}

\textbf{Mudanças climáticas e redução de risco de desastres:}

{
\begin{longtable}[]{@{}lcccl@{}}
\toprule\noalign{}
Impacto & Escala 7 pontos & Antes & Após & Comentários \\
\midrule\noalign{}
\endhead
\bottomrule\noalign{}
\endlastfoot
Impactos de enchentes (aum. $\leftarrow$ $\rightarrow$ dim.) & $\circ$$\circ$$\circ$$\circ$$\circ$$\circ$$\circ$ & & & \\
Deslizamentos/fluxos de detritos & $\circ$$\circ$$\circ$$\circ$$\circ$$\circ$$\circ$ & & & \\
Impactos de seca & $\circ$$\circ$$\circ$$\circ$$\circ$$\circ$$\circ$ & & & \\
Impactos de ciclones/tempestades & $\circ$$\circ$$\circ$$\circ$$\circ$$\circ$$\circ$ & & & \\
Emissões de GEE (aum. $\leftarrow$ $\rightarrow$ dim.) & $\circ$$\circ$$\circ$$\circ$$\circ$$\circ$$\circ$ & & & \\
Risco de incêndio (aum. $\leftarrow$ $\rightarrow$ dim.) & $\circ$$\circ$$\circ$$\circ$$\circ$$\circ$$\circ$ & & & \\
Velocidade do vento (aum. $\leftarrow$ $\rightarrow$ dim.) & $\circ$$\circ$$\circ$$\circ$$\circ$$\circ$$\circ$ & & & \\
Microclima & $\circ$$\circ$$\circ$$\circ$$\circ$$\circ$$\circ$ & & & \\
\end{longtable}
}

\subsubsection{\texorpdfstring{6.2 Impactos fora do local (\emph{off-site}) {[Consulta]} {[Campo]}}{6.2 Impactos fora do local (off-site) {[Consulta]} {[Campo]}}}\label{impactos-fora-do-local-off-site}

{
\begin{longtable}[]{@{}lcccl@{}}
\toprule\noalign{}
Impacto & Escala 7 pontos & Antes & Após & Comentários \\
\midrule\noalign{}
\endhead
\bottomrule\noalign{}
\endlastfoot
Disponibilidade hídrica (lençol, nascentes) & $\circ$$\circ$$\circ$$\circ$$\circ$$\circ$$\circ$ & & & \\
Estabilidade de vazões & $\circ$$\circ$$\circ$$\circ$$\circ$$\circ$$\circ$ & & & \\
Enchentes a jusante (aum. $\leftarrow$ $\rightarrow$ dim.) & $\circ$$\circ$$\circ$$\circ$$\circ$$\circ$$\circ$ & & & \\
Assoreamento a jusante (aum. $\leftarrow$ $\rightarrow$ dim.) & $\circ$$\circ$$\circ$$\circ$$\circ$$\circ$$\circ$ & & & \\
Poluição águas subterrâneas/rios (aum. $\leftarrow$ $\rightarrow$ dim.) & $\circ$$\circ$$\circ$$\circ$$\circ$$\circ$$\circ$ & & & \\
Capacidade de filtragem/tamponamento & $\circ$$\circ$$\circ$$\circ$$\circ$$\circ$$\circ$ & & & \\
Sedimentos transportados pelo vento (aum. $\leftarrow$ $\rightarrow$ dim.) & $\circ$$\circ$$\circ$$\circ$$\circ$$\circ$$\circ$ & & & \\
Danos a campos vizinhos (aum. $\leftarrow$ $\rightarrow$ dim.) & $\circ$$\circ$$\circ$$\circ$$\circ$$\circ$$\circ$ & & & \\
Danos a infraestrutura (aum. $\leftarrow$ $\rightarrow$ dim.) & $\circ$$\circ$$\circ$$\circ$$\circ$$\circ$$\circ$ & & & \\
Impacto de GEE (aum. $\leftarrow$ $\rightarrow$ dim.) & $\circ$$\circ$$\circ$$\circ$$\circ$$\circ$$\circ$ & & & \\
\end{longtable}
}

\subsubsection{6.3 Exposição e sensibilidade da Tecnologia a mudanças climáticas graduais e extremos/desastres {[Consulta]}}\label{exposiuxe7uxe3o-e-sensibilidade-da-tecnologia-a-mudanuxe7as-climuxe1ticas-graduais-e-extremosdesastres}

\emph{Indique mudanças graduais no clima e extremos/desastres climáticos observados pelos usuários nos últimos 10 anos.}

\subsubsection{Mudanças climáticas graduais}\label{mudanuxe7as-climuxe1ticas-graduais}

{
\begin{longtable}[]{@{}
  >{\raggedright\arraybackslash}p{(\linewidth - 6\tabcolsep) * \real{0.1406}}
  >{\centering\arraybackslash}p{(\linewidth - 6\tabcolsep) * \real{0.1562}}
  >{\centering\arraybackslash}p{(\linewidth - 6\tabcolsep) * \real{0.1562}}
  >{\centering\arraybackslash}p{(\linewidth - 6\tabcolsep) * \real{0.5469}}@{}}
\toprule\noalign{}
\begin{minipage}[b]{\linewidth}\raggedright
Mudança
\end{minipage} & \begin{minipage}[b]{\linewidth}\centering
Aumentou
\end{minipage} & \begin{minipage}[b]{\linewidth}\centering
Diminuiu
\end{minipage} & \begin{minipage}[b]{\linewidth}\centering
Capacidade da Tecnologia de lidar
\end{minipage} \\
\midrule\noalign{}
\endhead
\bottomrule\noalign{}
\endlastfoot
Temperatura anual & $\circ$ & $\circ$ & $\circ$ Muito precária $\circ$ Precária $\circ$ Moderada $\circ$ Bem $\circ$ Muito bem \\
Temperatura sazonal (estação: \_\_\_\_) & $\circ$ & $\circ$ & $\circ$ Muito precária $\circ$ Precária $\circ$ Moderada $\circ$ Bem $\circ$ Muito bem \\
Precipitação anual & $\circ$ & $\circ$ & $\circ$ Muito precária $\circ$ Precária $\circ$ Moderada $\circ$ Bem $\circ$ Muito bem \\
Precipitação sazonal (estação: \_\_\_\_) & $\circ$ & $\circ$ & $\circ$ Muito precária $\circ$ Precária $\circ$ Moderada $\circ$ Bem $\circ$ Muito bem \\
\end{longtable}
}

\subsubsection{Extremos climáticos (desastres)}\label{extremos-climuxe1ticos-desastres}

\textbf{Desastres meteorológicos:}

{
\begin{longtable}[]{@{}
  >{\raggedright\arraybackslash}p{(\linewidth - 2\tabcolsep) * \real{0.1875}}
  >{\centering\arraybackslash}p{(\linewidth - 2\tabcolsep) * \real{0.8125}}@{}}
\toprule\noalign{}
\begin{minipage}[b]{\linewidth}\raggedright
Tipo
\end{minipage} & \begin{minipage}[b]{\linewidth}\centering
Capacidade da Tecnologia
\end{minipage} \\
\midrule\noalign{}
\endhead
\bottomrule\noalign{}
\endlastfoot
$\Box$ Tempestade tropical (ciclone, tufão, furacão) & $\circ$ Muito precária $\circ$ Precária $\circ$ Moderada $\circ$ Bem $\circ$ Muito bem \\
$\Box$ Ciclone extratropical (tempestade de inverno) & $\circ$ Muito precária $\circ$ Precária $\circ$ Moderada $\circ$ Bem $\circ$ Muito bem \\
$\Box$ Chuva intensa local & $\circ$ Muito precária $\circ$ Precária $\circ$ Moderada $\circ$ Bem $\circ$ Muito bem \\
$\Box$ Trovoada local & $\circ$ Muito precária $\circ$ Precária $\circ$ Moderada $\circ$ Bem $\circ$ Muito bem \\
$\Box$ Granizo local & $\circ$ Muito precária $\circ$ Precária $\circ$ Moderada $\circ$ Bem $\circ$ Muito bem \\
$\Box$ Tempestade de neve local & $\circ$ Muito precária $\circ$ Precária $\circ$ Moderada $\circ$ Bem $\circ$ Muito bem \\
$\Box$ Tempestade de areia/poeira local & $\circ$ Muito precária $\circ$ Precária $\circ$ Moderada $\circ$ Bem $\circ$ Muito bem \\
$\Box$ Vendaval local & $\circ$ Muito precária $\circ$ Precária $\circ$ Moderada $\circ$ Bem $\circ$ Muito bem \\
$\Box$ Tornado local & $\circ$ Muito precária $\circ$ Precária $\circ$ Moderada $\circ$ Bem $\circ$ Muito bem \\
\end{longtable}
}

\textbf{Desastres climatológicos:}

{
\begin{longtable}[]{@{}lc@{}}
\toprule\noalign{}
Tipo & Capacidade da Tecnologia \\
\midrule\noalign{}
\endhead
\bottomrule\noalign{}
\endlastfoot
$\Box$ Onda de calor & $\circ$$\circ$$\circ$$\circ$$\circ$ \\
$\Box$ Onda de frio (geada) & $\circ$$\circ$$\circ$$\circ$$\circ$ \\
$\Box$ Condições extremas de inverno & $\circ$$\circ$$\circ$$\circ$$\circ$ \\
$\Box$ Seca & $\circ$$\circ$$\circ$$\circ$$\circ$ \\
$\Box$ Incêndio florestal & $\circ$$\circ$$\circ$$\circ$$\circ$ \\
$\Box$ Incêndio de campo (capim, arbustos) & $\circ$$\circ$$\circ$$\circ$$\circ$ \\
\end{longtable}
}

\textbf{Desastres hidrológicos:}

{
\begin{longtable}[]{@{}lc@{}}
\toprule\noalign{}
Tipo & Capacidade da Tecnologia \\
\midrule\noalign{}
\endhead
\bottomrule\noalign{}
\endlastfoot
$\Box$ Enchente fluvial & $\circ$$\circ$$\circ$$\circ$$\circ$ \\
$\Box$ Enxurrada (\emph{flash flood}) & $\circ$$\circ$$\circ$$\circ$$\circ$ \\
$\Box$ Maré de tempestade / enchente costeira & $\circ$$\circ$$\circ$$\circ$$\circ$ \\
$\Box$ Deslizamento / fluxo de detritos & $\circ$$\circ$$\circ$$\circ$$\circ$ \\
$\Box$ Avalanche & $\circ$$\circ$$\circ$$\circ$$\circ$ \\
\end{longtable}
}

\textbf{Desastres biológicos:}

{
\begin{longtable}[]{@{}lc@{}}
\toprule\noalign{}
Tipo & Capacidade da Tecnologia \\
\midrule\noalign{}
\endhead
\bottomrule\noalign{}
\endlastfoot
$\Box$ Epidemias (virais, bacterianas, fúngicas, parasitárias) & $\circ$$\circ$$\circ$$\circ$$\circ$ \\
$\Box$ Infestações de insetos/vermes (gafanhotos, lagartas, etc.) & $\circ$$\circ$$\circ$$\circ$$\circ$ \\
\end{longtable}
}

\textbf{Outras consequências climáticas:}

{
\begin{longtable}[]{@{}lc@{}}
\toprule\noalign{}
Consequência & Capacidade da Tecnologia \\
\midrule\noalign{}
\endhead
\bottomrule\noalign{}
\endlastfoot
$\Box$ Período de crescimento ampliado & $\circ$$\circ$$\circ$$\circ$$\circ$ \\
$\Box$ Período de crescimento reduzido & $\circ$$\circ$$\circ$$\circ$$\circ$ \\
$\Box$ Elevação do nível do mar & $\circ$$\circ$$\circ$$\circ$$\circ$ \\
$\Box$ Outro: \_\_\_\_\_\_\_\_ & $\circ$$\circ$$\circ$$\circ$$\circ$ \\
\end{longtable}
}

\subsubsection{6.4 Análise custo-benefício {[Consulta]}}\label{anuxe1lise-custo-benefuxedcio}

\textbf{Como os benefícios se comparam aos custos de implantação (perspectiva do usuário)?}

{
\begin{longtable}[]{@{}
  >{\raggedright\arraybackslash}p{(\linewidth - 14\tabcolsep) * \real{0.0870}}
  >{\centering\arraybackslash}p{(\linewidth - 14\tabcolsep) * \real{0.1304}}
  >{\centering\arraybackslash}p{(\linewidth - 14\tabcolsep) * \real{0.1304}}
  >{\centering\arraybackslash}p{(\linewidth - 14\tabcolsep) * \real{0.1304}}
  >{\centering\arraybackslash}p{(\linewidth - 14\tabcolsep) * \real{0.1304}}
  >{\centering\arraybackslash}p{(\linewidth - 14\tabcolsep) * \real{0.1304}}
  >{\centering\arraybackslash}p{(\linewidth - 14\tabcolsep) * \real{0.1304}}
  >{\centering\arraybackslash}p{(\linewidth - 14\tabcolsep) * \real{0.1304}}@{}}
\toprule\noalign{}
\begin{minipage}[b]{\linewidth}\raggedright
\end{minipage} & \begin{minipage}[b]{\linewidth}\centering
Muito negativo
\end{minipage} & \begin{minipage}[b]{\linewidth}\centering
Negativo
\end{minipage} & \begin{minipage}[b]{\linewidth}\centering
Lev. negativo
\end{minipage} & \begin{minipage}[b]{\linewidth}\centering
Neutro/equilibrado
\end{minipage} & \begin{minipage}[b]{\linewidth}\centering
Lev. positivo
\end{minipage} & \begin{minipage}[b]{\linewidth}\centering
Positivo
\end{minipage} & \begin{minipage}[b]{\linewidth}\centering
Muito positivo
\end{minipage} \\
\midrule\noalign{}
\endhead
\bottomrule\noalign{}
\endlastfoot
Retorno de curto prazo (1--3 anos) & $\circ$ & $\circ$ & $\circ$ & $\circ$ & $\circ$ & $\circ$ & $\circ$ \\
Retorno de longo prazo (\textgreater{} 10 anos) & $\circ$ & $\circ$ & $\circ$ & $\circ$ & $\circ$ & $\circ$ & $\circ$ \\
\end{longtable}
}

\textbf{Como os benefícios se comparam aos custos de manutenção (perspectiva do usuário)?}

{
\begin{longtable}[]{@{}
  >{\raggedright\arraybackslash}p{(\linewidth - 14\tabcolsep) * \real{0.0870}}
  >{\centering\arraybackslash}p{(\linewidth - 14\tabcolsep) * \real{0.1304}}
  >{\centering\arraybackslash}p{(\linewidth - 14\tabcolsep) * \real{0.1304}}
  >{\centering\arraybackslash}p{(\linewidth - 14\tabcolsep) * \real{0.1304}}
  >{\centering\arraybackslash}p{(\linewidth - 14\tabcolsep) * \real{0.1304}}
  >{\centering\arraybackslash}p{(\linewidth - 14\tabcolsep) * \real{0.1304}}
  >{\centering\arraybackslash}p{(\linewidth - 14\tabcolsep) * \real{0.1304}}
  >{\centering\arraybackslash}p{(\linewidth - 14\tabcolsep) * \real{0.1304}}@{}}
\toprule\noalign{}
\begin{minipage}[b]{\linewidth}\raggedright
\end{minipage} & \begin{minipage}[b]{\linewidth}\centering
Muito negativo
\end{minipage} & \begin{minipage}[b]{\linewidth}\centering
Negativo
\end{minipage} & \begin{minipage}[b]{\linewidth}\centering
Lev. negativo
\end{minipage} & \begin{minipage}[b]{\linewidth}\centering
Neutro/equilibrado
\end{minipage} & \begin{minipage}[b]{\linewidth}\centering
Lev. positivo
\end{minipage} & \begin{minipage}[b]{\linewidth}\centering
Positivo
\end{minipage} & \begin{minipage}[b]{\linewidth}\centering
Muito positivo
\end{minipage} \\
\midrule\noalign{}
\endhead
\bottomrule\noalign{}
\endlastfoot
Retorno de curto prazo & $\circ$ & $\circ$ & $\circ$ & $\circ$ & $\circ$ & $\circ$ & $\circ$ \\
Retorno de longo prazo & $\circ$ & $\circ$ & $\circ$ & $\circ$ & $\circ$ & $\circ$ & $\circ$ \\
\end{longtable}
}

Comentários: \_\_\_\_\_\_\_\_\_\_\_\_\_\_\_\_\_\_\_\_\_\_\_\_\_\_\_\_\_\_\_\_\_\_\_\_\_\_\_\_

\subsubsection{6.5 Adoção da Tecnologia {[Consulta]}}\label{adouxe7uxe3o-da-tecnologia}

Quantos usuários da terra na área adotaram/implementaram a Tecnologia?

$\circ$ Casos isolados/experimentais $\circ$ 1--10\% $\circ$ 10--50\% $\circ$ \textgreater{} 50\%

Se disponível, quantifique (nº de domicílios e/ou área coberta): \_\_\_\_\_\_\_\_\_\_\_\_\_\_\_\_\_\_\_\_\_\_\_\_\_\_\_\_\_\_\_\_\_\_\_\_\_\_\_\_

Dos que adotaram, quantos o fizeram espontaneamente (sem incentivos materiais/pagamentos)?

$\circ$ 0--10\% $\circ$ 10--50\% $\circ$ 50--90\% $\circ$ 90--100\%

Comentários: \_\_\_\_\_\_\_\_\_\_\_\_\_\_\_\_\_\_\_\_\_\_\_\_\_\_\_\_\_\_\_\_\_\_\_\_\_\_\_\_

\subsubsection{6.6 Adaptação {[Consulta]}}\label{adaptauxe7uxe3o}

A Tecnologia foi modificada recentemente para adaptar-se a condições em mudança?

$\circ$ Não $\circ$ Sim

Se sim, adaptada a:

$\circ$ Mudanças/extremos climáticos $\circ$ Mercados em transformação $\circ$ Disponibilidade de mão de obra (p.ex. migração) $\circ$ Outro: \_\_\_\_\_\_\_\_\_\_\_\_\_\_\_\_\_\_\_\_\_\_\_\_\_\_\_\_\_\_\_\_\_\_\_\_\_\_\_\_

Especifique a adaptação (projeto, material/espécies, etc.): \_\_\_\_\_\_\_\_\_\_\_\_\_\_\_\_\_\_\_\_\_\_\_\_\_\_\_\_\_\_\_\_\_\_\_\_\_\_\_\_

\subsubsection{6.7 Pontos fortes / vantagens / oportunidades da Tecnologia}\label{pontos-fortes-vantagens-oportunidades-da-tecnologia}

\emph{Diferencie entre as perspectivas de usuários da terra e pessoas-recurso.}

\textbf{Perspectiva do usuário da terra:}

\begin{enumerate}
\def\labelenumi{\arabic{enumi}.}
\item
  \begin{center}\rule{0.5\linewidth}{0.5pt}\end{center}
\item
  \begin{center}\rule{0.5\linewidth}{0.5pt}\end{center}
\item
  \begin{center}\rule{0.5\linewidth}{0.5pt}\end{center}
\item
  \begin{center}\rule{0.5\linewidth}{0.5pt}\end{center}
\end{enumerate}

\textbf{Perspectiva do compilador ou outra pessoa-recurso:}

\begin{enumerate}
\def\labelenumi{\arabic{enumi}.}
\item
  \begin{center}\rule{0.5\linewidth}{0.5pt}\end{center}
\item
  \begin{center}\rule{0.5\linewidth}{0.5pt}\end{center}
\item
  \begin{center}\rule{0.5\linewidth}{0.5pt}\end{center}
\item
  \begin{center}\rule{0.5\linewidth}{0.5pt}\end{center}
\end{enumerate}

\subsubsection{6.8 Pontos fracos / desvantagens / riscos da Tecnologia e formas de superá-los}\label{pontos-fracos-desvantagens-riscos-da-tecnologia-e-formas-de-superuxe1-los}

\textbf{Perspectiva do usuário da terra:}

{
\begin{longtable}[]{@{}ll@{}}
\toprule\noalign{}
Fraqueza / desvantagem / risco & Como superar? \\
\midrule\noalign{}
\endhead
\bottomrule\noalign{}
\endlastfoot
1. \_\_\_\_\_\_\_\_ & 1. \_\_\_\_\_\_\_\_ \\
2. \_\_\_\_\_\_\_\_ & 2. \_\_\_\_\_\_\_\_ \\
3. \_\_\_\_\_\_\_\_ & 3. \_\_\_\_\_\_\_\_ \\
4. \_\_\_\_\_\_\_\_ & 4. \_\_\_\_\_\_\_\_ \\
\end{longtable}
}

\textbf{Perspectiva do compilador ou outra pessoa-recurso:}

{
\begin{longtable}[]{@{}ll@{}}
\toprule\noalign{}
Fraqueza / desvantagem / risco & Como superar? \\
\midrule\noalign{}
\endhead
\bottomrule\noalign{}
\endlastfoot
1. \_\_\_\_\_\_\_\_ & 1. \_\_\_\_\_\_\_\_ \\
2. \_\_\_\_\_\_\_\_ & 2. \_\_\_\_\_\_\_\_ \\
3. \_\_\_\_\_\_\_\_ & 3. \_\_\_\_\_\_\_\_ \\
4. \_\_\_\_\_\_\_\_ & 4. \_\_\_\_\_\_\_\_ \\
\end{longtable}
}

\newpage

\begin{center}\rule{0.5\linewidth}{0.5pt}\end{center}

\subsection{7. Referências e links}\label{referuxeancias-e-links}

\subsubsection{7.1 Métodos / fontes de informação}\label{muxe9todos-fontes-de-informauxe7uxe3o}

\emph{Várias respostas possíveis.}

{
\begin{longtable}[]{@{}ll@{}}
\toprule\noalign{}
Método/fonte & Especifique (p.ex. nº de informantes) \\
\midrule\noalign{}
\endhead
\bottomrule\noalign{}
\endlastfoot
$\Box$ Visitas/levantamentos de campo & \_\_\_\_\_\_\_\_ \\
$\Box$ Entrevistas com usuários da terra & \_\_\_\_\_\_\_\_ \\
$\Box$ Entrevistas com especialistas em MST & \_\_\_\_\_\_\_\_ \\
$\Box$ Compilação de relatórios e documentação existente & \_\_\_\_\_\_\_\_ \\
$\Box$ Outro: \_\_\_\_\_\_\_\_ & \_\_\_\_\_\_\_\_ \\
\end{longtable}
}

Data da coleta de dados (em campo): \_\_\_\_\_\_\_\_\_\_\_\_\_\_\_\_\_\_\_\_\_\_\_\_\_\_\_\_\_\_\_\_\_\_\_\_\_\_\_\_

Comentários: \_\_\_\_\_\_\_\_\_\_\_\_\_\_\_\_\_\_\_\_\_\_\_\_\_\_\_\_\_\_\_\_\_\_\_\_\_\_\_\_

\subsubsection{7.2 Referências a publicações disponíveis}\label{referuxeancias-a-publicauxe7uxf5es-disponuxedveis}

\emph{Liste publicações relevantes (relatórios, manuais, materiais de treinamento, estudos de caso, etc.).}

{
\begin{longtable}[]{@{}lll@{}}
\toprule\noalign{}
Título, autor, ano, ISBN & Disponível em & Custo \\
\midrule\noalign{}
\endhead
\bottomrule\noalign{}
\endlastfoot
& & \\
& & \\
& & \\
\end{longtable}
}

\subsubsection{7.3 Links para informações relevantes disponíveis online}\label{links-para-informauxe7uxf5es-relevantes-disponuxedveis-online}

{
\begin{longtable}[]{@{}ll@{}}
\toprule\noalign{}
Título / descrição & URL \\
\midrule\noalign{}
\endhead
\bottomrule\noalign{}
\endlastfoot
& \\
& \\
& \\
\end{longtable}
}

\subsubsection{7.4 Comentários gerais}\label{comentuxe1rios-gerais}

\emph{(p.ex. feedback sobre o questionário ou banco de dados, ou observações gerais)}

\begin{center}\rule{0.5\linewidth}{0.5pt}\end{center}

\begin{center}\rule{0.5\linewidth}{0.5pt}\end{center}

\newpage

\begin{center}\rule{0.5\linewidth}{0.5pt}\end{center}

\subsection{9. Seção Suplementar --- Itens para Avaliação de Vulnerabilidade Biocultural ($V_1$--$V_6$)}\label{secao-suplementar-v1v6}

\emph{Esta seção contém itens suplementares não cobertos pelo questionário WOCAT original, desenvolvidos para operacionalizar as seis dimensões de vulnerabilidade biocultural ($V_1$--$V_6$) em contexto de comunidades quilombolas. Os itens estão organizados em três camadas de coleta conforme a natureza da variável.}

\subsubsection{Camada 1 --- Autorrelato (respondido pelo agricultor/líder comunitário)}

\paragraph{9.1 Transmissão intergeracional e dimensão ritual ($V_1$ --- Erosão Intergeracional)}~\\[4pt]

\textbf{9.1.1 Participação de jovens na prática}\\[2pt]
Qual a proporção de jovens ($<$\,35 anos) que dominam esta prática/saber na comunidade?\\
$\circ$~Nenhum \quad $\circ$~Poucos ($<$\,25\%) \quad $\circ$~Alguns (25--50\%) \quad $\circ$~A maioria (50--75\%) \quad $\circ$~Quase todos ($>$\,75\%)

\medskip\noindent\textbf{9.1.2 Taxa de transmissão ativa}\\[2pt]
Com que frequência ocorrem eventos de ensino/aprendizagem desta prática entre gerações (rodas de conversa, acompanhamento na roça, mutirões com aprendizes)?\\
$\circ$~Nunca \quad $\circ$~Raramente ($<$\,1 vez/ano) \quad $\circ$~Anualmente \quad $\circ$~Semestralmente \quad $\circ$~Mensalmente ou mais

\medskip\noindent\textbf{9.1.3 Modalidade de transmissão}\\[2pt]
Como este saber é transmitido? (Várias respostas possíveis.)\\
$\Box$~Oralidade (conversa, narrativa, cantos)\\
$\Box$~Prática acompanhada (aprendiz trabalha junto ao mestre)\\
$\Box$~Ritual ou cerimônia (bênção de sementes, festas de colheita, oferendas)\\
$\Box$~Registro escrito ou audiovisual produzido pela comunidade\\
$\Box$~Não está sendo transmitido atualmente\\
$\Box$~Outro: \dotfill

\medskip\noindent\textbf{9.1.4 Significado ritual ou cosmológico vinculado à prática}\\[2pt]
Esta prática possui dimensão espiritual, ritual ou cosmológica (p.\,ex.\@ bênçãos sobre sementes, plantio associado a ciclos lunares, proibições em datas sagradas, oferendas)?\\
$\circ$~Não \quad $\circ$~Sim, secundário (complementar à função técnica) \quad $\circ$~Sim, central (inseparável da execução da prática)\\[2pt]
Se sim, descreva brevemente: \dotfill\\[2pt]\dotfill

\medskip\noindent\textbf{9.1.5 Percepção de risco de perda}\\[2pt]
Na sua avaliação, qual o risco de este saber desaparecer na próxima geração?\\
$\circ$~Muito baixo \quad $\circ$~Baixo \quad $\circ$~Moderado \quad $\circ$~Alto \quad $\circ$~Muito alto\\[2pt]
Comentários: \dotfill

\paragraph{9.2 Governança comunitária e capital social ($V_6$ --- Organização Social)}~\\[4pt]

\textbf{9.2.1 Mestres de saberes}\\[2pt]
Quantas pessoas na comunidade são reconhecidas como guardiãs/mestres deste saber?\\
$\circ$~Nenhuma \quad $\circ$~1--2 \quad $\circ$~3--5 \quad $\circ$~6--10 \quad $\circ$~$>$\,10\\[2pt]
Faixa etária predominante dos mestres: $\Box$~Jovens ($<$\,35) \quad $\Box$~Meia-idade (35--59) \quad $\Box$~Idosos ($\geq$\,60)

\medskip\noindent\textbf{9.2.2 Eventos coletivos de transmissão e prática}\\[2pt]
Com que frequência ocorrem eventos coletivos relacionados a este saber (mutirões, trocas de sementes, feiras de agrobiodiversidade, rodas de prosa, encontros intercomunitários)?\\
$\circ$~Nunca \quad $\circ$~Raramente ($<$\,1 vez/ano) \quad $\circ$~Anualmente \quad $\circ$~Semestralmente \quad $\circ$~Mensalmente ou mais

\medskip\noindent\textbf{9.2.3 Redes de cooperação intercomunitárias}\\[2pt]
A comunidade mantém intercâmbio deste saber com outras comunidades quilombolas ou tradicionais?\\
$\circ$~Não \quad $\circ$~Sim, ocasionalmente (contato informal, visitas esporádicas) \quad $\circ$~Sim, regularmente (rede organizada, feiras, projetos conjuntos)\\[2pt]
Se sim, especifique comunidades ou redes: \dotfill

\medskip\noindent\textbf{9.2.4 Governança de gênero no saber}\\[2pt]
Quem participa das decisões sobre o manejo desta prática?\\
$\Box$~Predominantemente mulheres \quad $\Box$~Predominantemente homens \quad $\Box$~Ambos em proporção equilibrada\\
$\Box$~Conselhos de anciãos/anciãs \quad $\Box$~Assembleia comunitária aberta \quad $\Box$~Outro: \dotfill\\[4pt]
Quem detém o conhecimento principal desta prática?\\
$\Box$~Predominantemente mulheres \quad $\Box$~Predominantemente homens \quad $\Box$~Ambos em proporção equilibrada

\subsubsection{Camada 2 --- Inventário etnobotânico e ecológico (agricultor + pesquisador em campo)}

\paragraph{9.3 Agrobiodiversidade e etnotaxonomia ($V_2$ --- Complexidade Biocultural)}~\\[4pt]

\textbf{9.3.1 Inventário de variedades locais}\\[2pt]
\emph{Instrução ao compilador: realize listagem livre com o agricultor. Para cada cultura principal, registre todas as variedades locais (crioulas, tradicionais) cultivadas ou conhecidas pelo informante. Use nomes locais.}

{
\begin{longtable}[]{@{}p{3cm}p{3.5cm}p{3.5cm}p{2.5cm}@{}}
\toprule\noalign{}
Cultura & Variedade local & Origem/história & Uso principal \\
\midrule\noalign{}
\endhead
\bottomrule\noalign{}
\endlastfoot
 &  &  &  \\
 &  &  &  \\
 &  &  &  \\
 &  &  &  \\
\end{longtable}
}

Número total de variedades locais listadas: \dotfill

\medskip\noindent\textbf{9.3.2 Classificação êmica de solos}\\[2pt]
\emph{Instrução ao compilador: registre como o agricultor classifica os tipos de solo, usando a terminologia local.}

{
\begin{longtable}[]{@{}p{2.5cm}p{3cm}p{3cm}p{3.5cm}@{}}
\toprule\noalign{}
Nome local & Características (agricultor) & Culturas associadas & Corresp. pedológica (pesquisador) \\
\midrule\noalign{}
\endhead
\bottomrule\noalign{}
\endlastfoot
 &  &  &  \\
 &  &  &  \\
 &  &  &  \\
\end{longtable}
}

\medskip\noindent\textbf{9.3.3 Interações solo-planta-clima codificadas}\\[2pt]
O agricultor segue regras de associação entre tipo de solo, variedade cultivada e época de plantio baseadas em conhecimento tradicional?\\
$\circ$~Não \quad $\circ$~Sim\\[2pt]
Se sim, descreva até 3 regras principais:\\
1. \dotfill\\
2. \dotfill\\
3. \dotfill

\paragraph{9.4 Singularidade e exclusividade territorial ($V_3$ --- Singularidade Territorial)}~\\[4pt]

\textbf{9.4.1 Exclusividade da prática}\\[2pt]
Até onde o agricultor sabe, esta prática/variedade existe em outras comunidades da região?\\
$\circ$~Sim, é comum na região \quad $\circ$~Existe em algumas comunidades \quad $\circ$~Só existe nesta comunidade \quad $\circ$~Não sabe

\medskip\noindent\textbf{9.4.2 Especificidade ambiental}\\[2pt]
Esta prática depende de condições ambientais específicas deste território (tipo de solo, microclima, água, vegetação nativa)?\\
$\circ$~Não, pode ser aplicada em qualquer lugar \quad $\circ$~Parcialmente \quad $\circ$~Sim, totalmente (só funciona aqui)\\[2pt]
Se sim, especifique: \dotfill

\medskip\noindent\textbf{9.4.3 Inventário comparativo entre comunidades}\\[2pt]
\emph{Instrução ao compilador: este item exige levantamento multi-sítio. Preencha após documentar pelo menos 2 comunidades. Para cada prática/saber, indique presença/ausência por comunidade.}

{
\begin{longtable}[]{@{}p{3.5cm}cccc@{}}
\toprule\noalign{}
Prática/saber & Com. A & Com. B & Com. C & Com. D \\
\midrule\noalign{}
\endhead
\bottomrule\noalign{}
\endlastfoot
 & $\Box$\,P $\Box$\,A & $\Box$\,P $\Box$\,A & $\Box$\,P $\Box$\,A & $\Box$\,P $\Box$\,A \\
 & $\Box$\,P $\Box$\,A & $\Box$\,P $\Box$\,A & $\Box$\,P $\Box$\,A & $\Box$\,P $\Box$\,A \\
 & $\Box$\,P $\Box$\,A & $\Box$\,P $\Box$\,A & $\Box$\,P $\Box$\,A & $\Box$\,P $\Box$\,A \\
\end{longtable}
}

\subsubsection{Camada 3 --- Auditoria documental e jurídico-institucional (preenchido pelo pesquisador/compilador)}

\paragraph{9.5 Diagnóstico de déficit documental ($V_4$ --- Status de Documentação)}~\\[4pt]

\textbf{9.5.1 Registro formal prévio}\\[2pt]
Este saber/prática já foi objeto de algum registro formal antes deste levantamento?\\
$\circ$~Não \quad $\circ$~Sim, parcialmente \quad $\circ$~Sim, completamente\\[2pt]
Se sim, tipo de registro:\\
$\Box$~Publicação acadêmica (artigo, dissertação, tese)\\
$\Box$~Relatório técnico (EMATER, EMBRAPA, ONG)\\
$\Box$~Ficha WOCAT ou banco de dados internacional\\
$\Box$~Registro audiovisual (vídeo, áudio, fotografia documental)\\
$\Box$~Inventário patrimonial (IPHAN, secretaria de cultura)\\
$\Box$~Registro comunitário (caderno, cartilha, protocolo da associação)\\
$\Box$~Outro: \dotfill\\[2pt]
Referência (se disponível): \dotfill

\medskip\noindent\textbf{9.5.2 Completude documental}\\[2pt]
\emph{Instrução ao compilador: para o conjunto de saberes/práticas identificados nesta comunidade, avalie o grau de documentação existente.}

{
\begin{longtable}[]{@{}p{9cm}c@{}}
\toprule\noalign{}
Indicador & Valor \\
\midrule\noalign{}
\endhead
\bottomrule\noalign{}
\endlastfoot
Número total de saberes/práticas identificados nesta comunidade & \\
Número com algum tipo de registro formal & \\
Número com ficha técnica completa (descrição, fotos, especificações) & \\
Número com registro audiovisual & \\
Número com registro digitalizado e acessível online & \\
\end{longtable}
}

\medskip\noindent\textbf{9.5.3 Risco de perda documental}\\[2pt]
Existem acervos orais (gravações, relatos) em risco de perda por deterioração de mídias, falecimento de informantes ou falta de infraestrutura de armazenamento?\\
$\circ$~Não \quad $\circ$~Sim\\[2pt]
Se sim, descreva: \dotfill

\paragraph{9.6 Proteção jurídica e propriedade intelectual ($V_5$ --- Vulnerabilidade Jurídica)}~\\[4pt]

\textbf{9.6.1 Situação fundiária da comunidade}

{
\begin{longtable}[]{@{}p{6cm}p{6.5cm}@{}}
\toprule\noalign{}
Item & Resposta \\
\midrule\noalign{}
\endhead
\bottomrule\noalign{}
\endlastfoot
Certificação da Fundação Cultural Palmares & $\circ$~Sim (Ano: \_\_\_\_) \quad $\circ$~Não \quad $\circ$~Em processo \\
Título de propriedade coletiva (INCRA/ITERBA) & $\circ$~Sim (Ano: \_\_\_\_) \quad $\circ$~Não \quad $\circ$~Em processo \\
Existência de conflito fundiário ativo & $\circ$~Sim \quad $\circ$~Não \\
Área titulada (ha) & \_\_\_\_\_\_\_\_\_\_\_\_ \\
Área reivindicada (ha) & \_\_\_\_\_\_\_\_\_\_\_\_ \\
\end{longtable}
}

\medskip\noindent\textbf{9.6.2 Instrumentos de propriedade intelectual}\\[2pt]
A comunidade possui ou está em processo de obtenção de algum dos instrumentos abaixo?

{
\begin{longtable}[]{@{}p{5.5cm}p{7cm}@{}}
\toprule\noalign{}
Instrumento & Situação \\
\midrule\noalign{}
\endhead
\bottomrule\noalign{}
\endlastfoot
Indicação Geográfica (IG --- INPI) & $\circ$~Não \quad $\circ$~Em processo \quad $\circ$~Sim (nº: \_\_\_\_) \\
Marca Coletiva (INPI) & $\circ$~Não \quad $\circ$~Em processo \quad $\circ$~Sim (nº: \_\_\_\_) \\
Registro de Patrimônio Imaterial (IPHAN) & $\circ$~Não \quad $\circ$~Em processo \quad $\circ$~Sim (nº: \_\_\_\_) \\
Registro estadual/municipal de patrimônio & $\circ$~Não \quad $\circ$~Em processo \quad $\circ$~Sim (órgão: \_\_\_\_) \\
Protocolo Comunitário (Nagoya/Kunming-Montreal) & $\circ$~Não \quad $\circ$~Em elaboração \quad $\circ$~Sim \\
Registro de cultivar tradicional (MAPA/RNC) & $\circ$~Não \quad $\circ$~Em processo \quad $\circ$~Sim \\
Selo ou certificação & $\circ$~Não \quad $\circ$~Em processo \quad $\circ$~Sim (\_\_\_\_\_\_) \\
\end{longtable}
}

\medskip\noindent\textbf{9.6.3 Conhecimento sobre mecanismos de proteção}\\[2pt]
A comunidade tem conhecimento da existência dos instrumentos legais listados acima?\\
$\circ$~Desconhece totalmente \quad $\circ$~Conhece vagamente \quad $\circ$~Conhece, mas não sabe como acessar \quad $\circ$~Conhece e já tentou acessar\\[4pt]
A comunidade já recebeu assessoria jurídica ou técnica sobre proteção de conhecimentos tradicionais?\\
$\circ$~Não \quad $\circ$~Sim (instituição: \dotfill)

\medskip\noindent\textbf{9.6.4 Casos de apropriação indevida}\\[2pt]
Há registro ou percepção de que saberes, variedades ou práticas da comunidade tenham sido apropriados por terceiros sem consentimento ou repartição de benefícios?\\
$\circ$~Não \quad $\circ$~Sim \quad $\circ$~Suspeita, mas sem confirmação\\[2pt]
Se sim, descreva brevemente: \dotfill\\[2pt]\dotfill

\bigskip\noindent\emph{Fim da Seção Suplementar.}

\newpage

\begin{center}\rule{0.5\linewidth}{0.5pt}\end{center}

\subsection{8. Anexo --- Listas de referência}\label{anexo-listas-de-referuxeancia}

\emph{As listas abaixo são empregadas nos campos de seleção do questionário (seções 3.2 e 3.3). Apresenta-se versão condensada.}

\subsubsection{Culturas anuais (lista parcial WOCAT--IPCC)}\label{culturas-anuais-lista-parcial-wocatipcc}

{
\begin{longtable}[]{@{}
  >{\raggedright\arraybackslash}p{(\linewidth - 2\tabcolsep) * \real{0.4118}}
  >{\raggedright\arraybackslash}p{(\linewidth - 2\tabcolsep) * \real{0.5882}}@{}}
\toprule\noalign{}
\begin{minipage}[b]{\linewidth}\raggedright
Grupo
\end{minipage} & \begin{minipage}[b]{\linewidth}\raggedright
Exemplos
\end{minipage} \\
\midrule\noalign{}
\endhead
\bottomrule\noalign{}
\endlastfoot
Cereais & milho, arroz (várzea/sequeiro), sorgo, milheto, trigo, aveia, centeio, cevada, quinoa, amaranto \\
Leguminosas e oleaginosas & feijão, lentilha, soja, ervilha, amendoim, girassol, mamona, gergelim \\
Raízes e tubérculos & mandioca, batata, batata-doce, inhame, taro, beterraba-açucareira \\
Hortaliças & tomate, cebola, alho, abóbora, berinjela, folhosas (alface, couve, espinafre), cenoura, pepino \\
Fibras e flores & algodão, linho, cânhamo, roseiras, tulipas \\
Forrageiras & alfafa, trevo, gramíneas \\
Medicinais/aromáticas & diversas \\
Tabaco & tabaco \\
\end{longtable}
}

\subsubsection{Culturas perenes e arbóreo-arbustivas (lista parcial)}\label{culturas-perenes-e-arbuxf3reo-arbustivas-lista-parcial}

{
\begin{longtable}[]{@{}
  >{\raggedright\arraybackslash}p{(\linewidth - 2\tabcolsep) * \real{0.4118}}
  >{\raggedright\arraybackslash}p{(\linewidth - 2\tabcolsep) * \real{0.5882}}@{}}
\toprule\noalign{}
\begin{minipage}[b]{\linewidth}\raggedright
Grupo
\end{minipage} & \begin{minipage}[b]{\linewidth}\raggedright
Exemplos
\end{minipage} \\
\midrule\noalign{}
\endhead
\bottomrule\noalign{}
\endlastfoot
Perenes não lenhosas & cana-de-açúcar, banana/plátano, abacaxi, sisal, capins forrageiros, maracujá \\
Arbóreo-arbustivas & café (sol/sombra), cacau, coco, dendê, manga, abacate, citros, goiaba, mamão, seringueira, teca, mogno, oliveira, uva, romã, figo, tamarindo, alfarroba, argan, cajueiro, castanha-do-pará, pistache, nozes, amêndoas \\
\end{longtable}
}

\subsubsection{Pecuária e produtos}\label{pecuuxe1ria-e-produtos}

{
\begin{longtable}[]{@{}
  >{\raggedright\arraybackslash}p{(\linewidth - 2\tabcolsep) * \real{0.2400}}
  >{\raggedright\arraybackslash}p{(\linewidth - 2\tabcolsep) * \real{0.7600}}@{}}
\toprule\noalign{}
\begin{minipage}[b]{\linewidth}\raggedright
Tipo
\end{minipage} & \begin{minipage}[b]{\linewidth}\raggedright
Produtos/serviços
\end{minipage} \\
\midrule\noalign{}
\endhead
\bottomrule\noalign{}
\endlastfoot
Bovinos (leite, corte, trabalho) & carne, leite, couro, tração, esterco \\
Bubalinos, equinos, muares, asininos & transporte, tração \\
Suínos & carne \\
Caprinos, ovinos & carne, leite, lã, couro \\
Aves, coelhos & carne, ovos \\
Camelos, dromedários & transporte, leite, carne \\
Apicultura & mel, cera, pólen \\
Piscicultura & peixes \\
Fauna silvestre & herbívoros grandes/pequenos \\
\end{longtable}
}

\subsubsection{Tipos de florestas (lista parcial)}\label{tipos-de-florestas-lista-parcial}

{
\begin{longtable}[]{@{}
  >{\raggedright\arraybackslash}p{(\linewidth - 2\tabcolsep) * \real{0.3667}}
  >{\raggedright\arraybackslash}p{(\linewidth - 2\tabcolsep) * \real{0.6333}}@{}}
\toprule\noalign{}
\begin{minipage}[b]{\linewidth}\raggedright
Categoria
\end{minipage} & \begin{minipage}[b]{\linewidth}\raggedright
Exemplos de zonas
\end{minipage} \\
\midrule\noalign{}
\endhead
\bottomrule\noalign{}
\endlastfoot
\textbf{Florestas naturais} & boreal conífera, boreal montana, tundra boreal, subtropical desértica/seca/úmida/montana, temperada continental/oceânica/desértica/montana/estepe, tropical seca/úmida decídua/montana/pluvial/arbustiva \\
\textbf{Plantações} & Por zona climática × gênero principal (Eucalyptus spp., Pinus spp., Tectona grandis, folhosas diversas) \\
\end{longtable}
}

\subsubsection{Tipos de árvores (lista parcial)}\label{tipos-de-uxe1rvores-lista-parcial}

Acacia spp., Eucalyptus spp., Pinus spp., Tectona spp., Bamboo spp., Casuarina spp., Grevillea robusta, Azadirachta indica, Populus spp., Salix spp., Leucaena leucocephala, Dalbergia sissoo, Hevea brasiliensis, Swietenia macrophylla, entre outros.

\begin{center}\rule{0.5\linewidth}{0.5pt}\end{center}

\emph{Fim do questionário.}


\end{document}
